\documentclass[10pt, a4paper]{report}
% IEEE body text = 10pt. Using report class for chapter support.

% =============================================================================
% PACKAGES
% =============================================================================

% --- Encoding & Language ---
\usepackage[utf8]{inputenc}
\usepackage[T1]{fontenc}
\usepackage[british]{babel}

% --- Typography ---
% IEEE uses Times Roman throughout
\usepackage{mathptmx}              % Times Roman for text + math
\usepackage[scaled=0.90]{helvet}   % Helvetica (sans-serif fallback)
\usepackage{courier}               % Courier for monospace / code

% --- Page Geometry (IEEE A4 spec) ---
% IEEE A4: side = 13mm, top = 19mm, bottom = 43mm
% Single-column width = 184mm
\usepackage[
    a4paper,
    left=13mm,
    right=13mm,
    top=19mm,
    bottom=43mm,
    headheight=12pt,
    headsep=10pt,
    footskip=20pt
]{geometry}

% --- Line Spacing ---
% IEEE body text: 10pt font / 12pt line spacing = factor of 1.2
% LaTeX default for 10pt is ~12pt baselineskip, so single spacing is correct.
\usepackage{setspace}
\singlespacing

% --- Headers & Footers ---
\usepackage{fancyhdr}
\pagestyle{fancy}
\fancyhf{}
\fancyhead[L]{\footnotesize\nouppercase{\leftmark}}
\fancyhead[R]{\footnotesize\thepage}
\renewcommand{\headrulewidth}{0.4pt}
\fancypagestyle{plain}{%
    \fancyhf{}
    \fancyfoot[C]{\footnotesize\thepage}
    \renewcommand{\headrulewidth}{0pt}
}

% --- Graphics & Figures ---
\usepackage{graphicx}
\usepackage{float}
\usepackage{subcaption}

% IEEE Figure captions: 8pt, 9pt line spacing, centred below figure
% IEEE Table captions: 8pt, 9pt line spacing, centred above table, small caps label
\usepackage[
    font={footnotesize, singlespacing},  % 8pt at 10pt base
    labelfont=bf,
    labelsep=period,
    justification=centering,
    tableposition=top,
    figureposition=bottom
]{caption}

% Override table caption label to use small caps "TABLE" with Roman numerals
\DeclareCaptionLabelFormat{ieeetable}{\textsc{Table~#2}}
\captionsetup[table]{labelformat=ieeetable, labelsep=newline, justification=centering}
\renewcommand{\thetable}{\Roman{table}}

% Figure numbering: Arabic, label "Fig."
\DeclareCaptionLabelFormat{ieeefigure}{Fig.~#2}
\captionsetup[figure]{labelformat=ieeefigure, labelsep=period, justification=centering}

% --- Tables ---
\usepackage{booktabs}
\usepackage{array}
\usepackage{multirow}
\usepackage{tabularx}

% --- Mathematics ---
\usepackage{amsmath}
\usepackage{amssymb}
\usepackage{bm}
\usepackage{braket}

% --- Colours ---
\usepackage{xcolor}
\definecolor{codeBackground}{RGB}{248,248,250}
\definecolor{codeFrame}{RGB}{190,190,200}
\definecolor{codeKeyword}{RGB}{0,80,160}
\definecolor{codeComment}{RGB}{80,128,80}
\definecolor{codeString}{RGB}{160,48,48}
\definecolor{codeNumber}{RGB}{140,140,150}
\definecolor{linkBlue}{RGB}{0,60,120}

% --- Code Listings ---
\usepackage{listings}
\lstdefinestyle{ieeeCode}{
    backgroundcolor=\color{codeBackground},
    frame=single,
    rulecolor=\color{codeFrame},
    framesep=3pt,
    xleftmargin=6pt,
    xrightmargin=4pt,
    % IEEE code should be small --- 8pt matches footnote/caption size
    basicstyle=\ttfamily\scriptsize\linespread{1.0}\selectfont,
    keywordstyle=\color{codeKeyword}\bfseries,
    commentstyle=\color{codeComment}\itshape,
    stringstyle=\color{codeString},
    numberstyle=\tiny\color{codeNumber}\ttfamily,
    numbers=left,
    numbersep=6pt,
    breaklines=true,
    breakatwhitespace=false,
    showstringspaces=false,
    captionpos=b,
    tabsize=4,
    columns=flexible,
    keepspaces=true,
    aboveskip=10pt,
    belowskip=6pt,
    abovecaptionskip=4pt,
    belowcaptionskip=2pt,
    literate={-}{-}1
}
\lstset{style=ieeeCode}
\renewcommand{\lstlistingname}{Listing}
\renewcommand{\lstlistlistingname}{List of Listings}

% --- Algorithms ---
\usepackage[ruled, vlined, linesnumbered]{algorithm2e}
\SetAlgoLined
\SetKwComment{Comment}{$\triangleright$\ }{}
\SetAlCapSkip{6pt}
% Match algorithm caption size to IEEE caption spec (8pt)
\SetAlCapFnt{\footnotesize}
\SetAlCapNameFnt{\footnotesize\bfseries}

% --- Hyperlinks ---
\usepackage[
    colorlinks=true,
    linkcolor=black,
    citecolor=linkBlue,
    urlcolor=linkBlue,
    bookmarks=true,
    bookmarksnumbered=true,
    pdfauthor={Om Kalyan},
    pdftitle={Beyond Classical Cryptography: Simulating and Benchmarking QKD},
    pdfsubject={BEng Dissertation},
    pdfkeywords={QKD, BB84, B92, E91, Qiskit, Post-Quantum Cryptography}
]{hyperref}
\usepackage[noabbrev, capitalise]{cleveref}

% --- Bibliography ---
% IEEE references: 8pt, 9pt line spacing, bracketed Arabic numerals
\usepackage{cite}

% --- Lists ---
\usepackage{enumitem}
\setlist{nosep, leftmargin=*}

% --- Paragraph formatting ---
% IEEE uses indented paragraphs with no extra inter-paragraph space
\setlength{\parindent}{1em}
\setlength{\parskip}{0pt}

% =============================================================================
% HEADING STYLES --- IEEE Conference Manuscript Format Definition
% =============================================================================
\usepackage{titlesec}

% --- CHAPTER ─────────────────────────────────────────────────────────────────
% Dissertation addition (IEEE conf. papers have no chapters).
% Styled as a prominent centred heading, consistent with IEEE aesthetic.
% Uses 14pt bold small caps, centred, with "CHAPTER I" label above.
\titleformat{\chapter}[display]
    {\normalfont\fontsize{14}{16}\bfseries\scshape\centering}
    {\MakeUppercase{\chaptertitlename}~\thechapter}
    {6pt}
    {\fontsize{14}{16}\selectfont\bfseries}
\titlespacing*{\chapter}{0pt}{-10pt}{20pt}

% --- SECTION (IEEE "Section Heads (1)") ──────────────────────────────────────
% Spec: 10pt regular, 14pt line spacing, centred,
%       Roman numerals (I., II., ...), CAPITALS / SMALL CAPS
\renewcommand{\thesection}{\Roman{section}}
\titleformat{\section}
    {\normalfont\fontsize{10}{14}\selectfont\scshape\centering}
    {\thesection.}{0.5em}{}
\titlespacing*{\section}{0pt}{14pt}{8pt}

% --- SUBSECTION (IEEE "Heading 2") ──────────────────────────────────────────
% Spec: 10pt italic, 12pt line spacing, left-justified,
%       Capital alpha numerals (A., B., ...)
\renewcommand{\thesubsection}{\Alph{subsection}}
\titleformat{\subsection}
    {\normalfont\fontsize{10}{12}\selectfont\itshape}
    {\thesubsection.}{0.5em}{}
\titlespacing*{\subsection}{0pt}{10pt}{4pt}

% --- SUBSUBSECTION (IEEE "Heading 3") ───────────────────────────────────────
% Spec: 10pt italic, 12pt line spacing, left-justified,
%       Arabic numeral + right paren, run-in, indent 1em
\renewcommand{\thesubsubsection}{\arabic{subsubsection}}
\titleformat{\subsubsection}[runin]
    {\normalfont\fontsize{10}{12}\selectfont\itshape}
    {\thesubsubsection)}{0.5em}{}[:]
\titlespacing*{\subsubsection}{1em}{8pt}{0.5em}

% --- PARAGRAPH (IEEE "Heading 4") ───────────────────────────────────────────
% Spec: 10pt italic, left-justified,
%       Lowercase alpha + right paren + colon, run-in, indent 2em
\renewcommand{\theparagraph}{\alph{paragraph}}
\titleformat{\paragraph}[runin]
    {\normalfont\fontsize{10}{12}\selectfont\itshape}
    {\theparagraph)}{0.5em}{}[:]
\titlespacing*{\paragraph}{2em}{6pt}{0.5em}

% Enable numbering depth through paragraph level
\setcounter{secnumdepth}{4}
\setcounter{tocdepth}{3}

% =============================================================================
% CUSTOM COMMANDS (QKD Notation)
% =============================================================================

\newcommand{\QBER}{\ensuremath{\mathrm{QBER}}}
\newcommand{\keyrate}{\ensuremath{R_{\mathrm{sec}}}}

% =============================================================================
% FOOTNOTE SIZE (IEEE: 8pt, 9pt line spacing)
% =============================================================================
\renewcommand{\footnotesize}{\fontsize{8}{9}\selectfont}

% =============================================================================
% DOCUMENT
% =============================================================================
\begin{document}

% ─────────────────────────────────────────────────────────────────────────────
% TITLE PAGE
% ─────────────────────────────────────────────────────────────────────────────
\begin{titlepage}
    \centering
    \thispagestyle{empty}
    \vspace*{3cm}

    % --- Title: IEEE spec = 24pt regular, 28pt line spacing, centred ---
    {\fontsize{28}{32}\selectfont
    Beyond Classical Cryptography:\par
    Simulating \& Benchmarking\par
    Quantum Key Distribution\par}

    \vspace{1.0cm}

    % --- Author Name: IEEE spec = 11pt regular, 12pt line spacing, centred ---
    {\fontsize{14}{15}\selectfont
    Om Kalyan\par}

    \vspace{0.5cm}

    {\fontsize{10}{11}\selectfont
    230733440\par}

    \vspace{1.0cm}

    % --- Author Affiliation: IEEE spec = 10pt regular, 11pt line spacing, centred ---
    {\fontsize{10}{11}\selectfont
    A dissertation submitted in partial fulfilment of the requirements\\
    for the degree of\par}

    \vspace{0.4cm}

    {\fontsize{10}{11}\selectfont\bfseries
    Bachelor of Engineering\\
    Electronics \& Computer Engineering\par}

    \vspace{1.0cm}

    % --- University Crest ---
    
\includegraphics[width=0.24\textwidth]{misc images/ncl logo vector.png}

    \vspace{1.0cm}

    % --- Institution: 11pt to match author name tier ---
    {\fontsize{11}{13}\selectfont
    School of Engineering\\
    Newcastle University\par}

    \vspace{0.6cm}

    % --- Date ---
    {\fontsize{10}{11}\selectfont 2026\par}

\end{titlepage}

% ─────────────────────────────────────────────────────────────────────────────
% FRONT MATTER
% ─────────────────────────────────────────────────────────────────────────────
\pagenumbering{roman}

% --- Abstract ---
\clearpage
\chapter*{Abstract}
\addcontentsline{toc}{chapter}{Abstract}

% IEEE Abstract: 9pt bold, "Abstract---" bold italic, 10pt line spacing
{\fontsize{9}{10}\selectfont\bfseries
Quantum algorithms such as Shor's factorisation threaten public-key cryptography systems like RSA and ECC, which rely on computational hardness for security. Quantum Key Distribution (QKD) offers a fundamentally different security paradigm: rather than relying on computational complexity, QKD derives information-theoretic security from the laws of quantum mechanics, ensuring privacy regardless of adversarial computing power. This project presents a simulation and benchmarking platform implementing three QKD protocols: BB84, B92, and E91 in Python3 using the Qiskit SDK. The platform compares Prepare \& Measure approaches against Entanglement-Based key distribution, evaluating protocol performance first under ideal channel conditions to establish theoretical baselines, then under realistic constraints including eavesdropper presence, quantum bit error rates, and channel losses.\par
}

% --- Acknowledgements ---
\clearpage
\chapter*{Acknowledgements}
\addcontentsline{toc}{chapter}{Acknowledgements}

I would like to thank my supervisors Martin Johnston and Oktay \c{C}etinkaya for their guidance\ldots

% --- Table of Contents ---
\clearpage
\tableofcontents

% --- List of Figures ---
\clearpage
\listoffigures
\addcontentsline{toc}{chapter}{List of Figures}

% --- List of Tables ---
\clearpage
\listoftables
\addcontentsline{toc}{chapter}{List of Tables}

% --- List of Listings ---
\clearpage
\lstlistoflistings
\addcontentsline{toc}{chapter}{List of Listings}

% --- List of Algorithms ---
\clearpage
\listofalgorithms
\addcontentsline{toc}{chapter}{List of Algorithms}

% ─────────────────────────────────────────────────────────────────────────────
% MAIN CONTENT
% ─────────────────────────────────────────────────────────────────────────────
\clearpage
\pagenumbering{arabic}

% ═══════════════════════════════════════════════════════════════════════════════
\chapter{Introduction}
% ═══════════════════════════════════════════════════════════════════════════════

\section{A Brief History of Cryptology}

% TODO: A brief introduction to QKD --- why it is needed, who benefits
% from the technology, and where it is deployed in the real world.

% this should talk abput classical cryptography first. an overview of why we need secure systems
Since the dawn of days, humans have sought ways to communicate in secret. It only takes a glance back at the ancient societies of Mesopotamia, India, and China to understand where the foundations of cryptology were laid some 3000 years ago \cite{Kahn1967Code}. Whether it be the early use of Substitution or Transposition schemes, these early cryptographic systems established the core principle that would dominate cryptology for millennia: security through the obscurity of the method itself. 

Back in the 9th century, humanity experienced the first major analytical breakthrough when the Arab polymath Al-Kindi developed frequency analysis, a technique of cryptanalysis capable of breaking any monoalphabetic substitution cipher of the age by exploiting the statistical patterns of language itself \cite{AlKadi1992}. This initiated a centuries-long cycle of invention and defeat that defined the field of cryptology, with newly designed ciphers broken by increasingly sophisticated cryptanalysis. Fast-forward to the 16th century, Blaise de Vigen\`ere built on the work of Giovan Bellaso to introduce a polyalphabetic cipher that would go on to secure the world for 300 years, until the great minds of Babbage and Kasiski developed statistical methods to compromise the scheme in the 19th century \cite{Singh1999, Bauer2021}.

The dawn of the Second World War transformed cryptology from a niche discipline into a strategic imperative. Repeating this "Tug of War", Germany's Engima, a rotor-based electromechanical cipher, was famously broken at Bletchley Park by the Father of Theoretical Computer Science - Alan Turing, and his colleagues \cite{Hodges2014}. Despite this technological revolution, security through simplicity was demonstrated in the Pacific theatre by the US Marine Corps' deployment of Navajo Code Talkers who used their native language as the basis for a battlefield cipher. A truly innovative system that rested security on the extreme phonetic distance between the Navajo and Japanese languages remained unbroken until its declassification in 1968 \cite{Huffman2000}. 

Looking to the post-war era, Claude Shannon redefined cryptography using rigorous mathematical foundations in his initially classified 1945 paper on secrecy systems, introducing concepts such as confusion, diffusion, and perfect secrecy \cite{Shannon1945}. This revolutionary framework, later publicly published in 1948, fueled the development of the modern cryptographic infrastructure we rely on today: the Data Encryption Standard (DES) in 1977, the RSA public-key cryptosystem in 1978, and the Advanced Encryption Standard (AES) in 2001 \cite{Shannon1948, DiffieHellman1976, RSA1978}. Peeling back the layers of essentially all the digital communications we use today, from online banking and e-commerce to government and military networks, we find these algorithms underpinning the security on which we rely. 

\section{The Quantum Threat to Classical Cryptography}

% TODO: Shor's algorithm, the quantum computing roadmap, and the need for
% information-theoretically secure key exchange.
% new intro para
The security of modern public-key cryptography rests on the computational hardness of problems like integer factorisation and the discrete logarithm. Shor's algorithm \cite{Shor1997} demonstrated that a sufficiently powerful quantum computer could solve both in polynomial time, rendering schemes such as RSA and ECC fundamentally insecure. Two complementary responses have emerged: Post-Quantum Cryptography (PQC), which replaces vulnerable algorithms with quantum-resistant mathematical problems, and Quantum Key Distribution (QKD), which exploits the laws of quantum mechanics to distribute encryption keys with information-theoretic security. This project focuses on the latter, building a simulation and benchmarking platform for QKD protocols including BB84, B92, and E91.


\newpage
\section{Aims \& Objectives}

% TODO: Outline the project aims and objectives.
% Describe the overall research message and trajectory.
The aim of this project is to design, implement, and evaluate a modular simulation and benchmarking platform for Discrete-Variable Quantum Key Distribution protocols, enabling quantitative comparison of protocol performance under realistic channel conditions, eavesdropping attacks, and post-processing constraints. All simulations are conducted in Python using IBM's Qiskit framework and the Aer noise simulator. The specific objectives are organised into six categories.

\subsection{Literature Review \& Analytical Framework}

\begin{enumerate}[label=\textbf{O1.\arabic*}, leftmargin=1.5cm]
    \item Conduct a critical review of QKD protocol literature, covering the theoretical foundations of BB84 \cite{BB84}, B92 \cite{B92}, and E91 \cite{E91}, their respective security proofs, and known practical limitations. This review will establish the analytical framework and parameter choices (noise models, error correction efficiency, security bounds) used throughout the project.

    \item Identify and document the key metrics by which QKD protocols are evaluated in the literature, such as: Quantum Bit Error Rate (QBER), sifted key rate, secure key rate, maximum tolerable error threshold, and sifting efficiency. These metrics will form the basis of all subsequent benchmarking.
\end{enumerate}

\subsection{Protocol Implementation}

\begin{enumerate}[label=\textbf{O2.\arabic*}, leftmargin=1.5cm]
    \item \textbf{BB84 (Four-State, Two-Basis):} Implement the BB84 protocol as the primary baseline simulation. The implementation models Alice's random bit and basis selection, qubit state preparation in the rectilinear ($Z$) or diagonal ($X$) basis, quantum transmission through a configurable noisy channel, Bob's random basis measurement, and public basis reconciliation (sifting). The implementation is considered complete when the sifting efficiency converges to the theoretically expected $\sim$50\% across repeated simulation runs. % refine with simualtions

    \item \textbf{B92 (Two Non-Orthogonal States):} Implement Bennett's 1992 protocol as the primary comparative prepare and measure protocol, encoding bit values using only two non-orthogonal states ($|0\rangle$ and $|+\rangle$). Bob performs a measurement that can yield a conclusive result, an inconclusive result, or a photon loss event. The implementation is validated by confirming that the sifting efficiency converges to $\sim$25\% and that the QBER threshold falls below that of BB84, consistent with B92's reduced noise tolerance.

    \item \textbf{E91 (Entanglement-Based):} Implement Ekert's 1991 protocol. This models the generation and distribution of maximally entangled Bell pairs, independent random basis measurements from three orientations per party, key extraction from matching-basis outcomes, and eavesdropper detection via the CHSH inequality. The implementation is validated by confirming that the computed CHSH parameter $|S|$ converges to $2\sqrt{2} \approx 2.828$ in the absence of eavesdropping, and degrades towards the classical bound of 2 under increasing interception rate.
\end{enumerate}

\subsection{Channel Modelling and Eavesdropper Simulation}

\begin{enumerate}[label=\textbf{O3.\arabic*}, leftmargin=1.5cm]
    \item Model quantum channel effects using Qiskit Aer's noise framework: depolarising noise, and phase damping. Noise is applied exclusively to the quantum channel via an identity gate channel marker, not to local preparation or measurement operations, ensuring the simulation accurately represents a protocol where Alice's and Bob's local devices are trusted.
    
    \item Implement an intercept-resend eavesdropping attack with a configurable interception probability $p \in [0,1]$. Eve intercepts each qubit with probability $p$, measures in a randomly chosen basis, and re-transmits to Bob. Sweep the interception rate from 0\% to 100\%  record the resulting QBER for each protocol.
    
\end{enumerate}

\subsection{Post-Processing and Secure Key Rate Analysis}

\begin{enumerate}[label=\textbf{O4.\arabic*}, leftmargin=1.5cm]
    \item Calculate the secure key rate using the GLLP formula \cite{GLLP2004}, incorporating the error correction efficiency factor $f_{EC} = 1.16$ corresponding to the CASCADE protocol \cite{Cascade1994}. This yields a practical BB84 QBER security threshold of $\sim$9.8\%, compared to the idealised 11\% Shor-Preskill bound at $f_{EC} = 1.0$ \cite{ShorPreskill2000}. Document and justify this distinction, as it has direct engineering significance for determining whether a given fibre link can sustain secure key generation.

    \item Compute the mutual information quantities $I(A{:}B)$ (Alice-Bob) and $I(A{:}E)$ (Alice-Eve) as functions of QBER for both BB84 and B92. Identify the crossover point where $I(A{:}E) \geq I(A{:}B)$, beyond which no secret key can be extracted, and confirm this aligns with the GLLP-derived thresholds for each protocol.

    \item Determine and compare the maximum tolerable QBER thresholds for BB84 ($\sim$9.8\% with $f_{EC} = 1.16$) and B92 ($\sim$6.5\%), quantifying the practical noise tolerance gap between the two prepare and measure protocols.
\end{enumerate}

\subsection{Verification, Validation, and Benchmarking Visualisation}

\begin{enumerate}[label=\textbf{O5.\arabic*}, leftmargin=1.5cm]
    \item \textbf{Verification against theoretical predictions.} Confirm that simulation outputs match known analytical results:
    \begin{itemize}
        \item BB84 sifting efficiency converges to 50\% ($\pm$2\%).
        \item Full intercept-resend attack on BB84 produces QBER of 25\% ($\pm$2\%).
        \item Secure key rate reaches zero at the GLLP-predicted QBER thresholds for each protocol.
        \item B92 sifting efficiency converges to 25\%.
    \end{itemize}
    These checks serve as pass/fail criteria confirming that the simulation correctly implements the underlying quantum mechanics and post-processing mathematics.

    % potential edit on where i may want to validate results
    \item \textbf{Validation against published experimental data.} Compare simulated key rate versus distance curves against published experimental QKD demonstrations, including long-distance BB84 results \cite{Boaron2018} and metropolitan network deployments reported in recent survey literature \cite{Zhang2025}. The simulation is not expected to reproduce experimental results exactly (real systems include imperfections beyond the scope of this model), but the trends, orders of magnitude, and maximum operational distances should be broadly consistent. Discrepancies are documented and attributed to specific modelling assumptions.

    \item \textbf{Benchmarking visualisation suite.} Develop a set of figures to present simulation results and protocol comparisons, including:
    \begin{itemize}
        \item QBER versus eavesdropper interception rate for each protocol.
        \item Secure key rate versus fibre distance under realistic attenuation and $f_{EC} = 1.16$.
        \item Noise model comparison showing QBER response across depolarising, bit-flip, amplitude damping, and phase damping channels at varying strengths.
        \item Mutual information crossover plot ($I(A{:}B)$ vs.\ $I(A{:}E)$ vs.\ QBER).
    \end{itemize}
\end{enumerate}

\subsection{Real-World Feasibility and Broader Context}

\begin{enumerate}[label=\textbf{O6.\arabic*}, leftmargin=1.5cm]
    \item Evaluate the practical viability of BB84 and B92 over metropolitan fibre links (10-50 km), identifying the distance and noise regimes in which each protocol sustains a positive secure key rate with $f_{EC} = 1.16$. Determine, for a given noise floor, the maximum deployment distance at which each protocol remains operational.

    \item Contextualise the simulation results within the broader landscape of quantum-safe security, including the complementary roles of QKD and post-quantum cryptography (NIST FIPS 203/204/205), the "harvest now, decrypt later" threat model motivating near-term QKD deployment, and the economic and infrastructure barriers to adoption across sectors such as Research \& Development, Government, Banking \& Finance, Critical Infrastructure, Healthcare, and Emerging Industries.
\end{enumerate}

% ═══════════════════════════════════════════════════════════════════════════════    
\chapter{Literature Review}
% ═══════════════════════════════════════════════════════════════════════════════

% TODO: Review of QKD theory, protocol evolution, real-world deployments,
% and positioning of this work relative to existing research.

% repurpose the intro section into this paragraph

\section{The Quantum Threat to Classical Cryptography}
% new
The security of public-key cryptosystems such as RSA, Diffie-Hellman, and ECC depends on the assumed intractability of integer factorisation and the discrete logarithm problem. Shor's algorithm \cite{Shor1997} undermines both assumptions, solving each in polynomial time on a quantum computer. While no machine of sufficient scale exists today, resource estimates have fallen sharply. Gidney and Eker\r{a}~\cite{Gidney2021} showed that a 2048-bit RSA modulus could be factored in approximately eight hours using 20 million noisy qubits, and a 2025 follow-up by Gidney \cite{Gidney2025} reduced the requirement to under one million noisy qubits at the cost of roughly one week of runtime. As illustrated in Fig. \ref{fig:quantumAdvantage}, these declining estimates suggest we are approaching Quantum Economic Advantage (QEA): the point at which the practical resource cost of such an attack becomes feasible.

\begin{figure}[h]
    \centering
    \includegraphics[width=1\linewidth]{images.lit/quantumAdvantage.png}
    \caption{The Quantum Economic Advantage}
    \label{fig:quantumAdvantage}
\end{figure}

Symmetric-key cryptography is not immune. Grover's search algorithm \cite{Grover1996} provides a quadratic speed-up over brute-force key search, effectively halving the security level of schemes such as AES-128 and AES-256 \cite{QKDbible}. In practice, this means AES-128 offers only 64 bits of security against a quantum adversary, making AES-256 the minimum acceptable key length in a post-quantum setting. The asymmetry of impact is notable, where Shor's algorithm breaks public-key schemes outright, Grover's merely weakens symmetric schemes to a degree that can be compensated by doubling key sizes. It is for this reason that the quantum threat is primarily directed at public-key infrastructure.

The "harvest now, decrypt later" attack model gives these theoretical vulnerabilities immediate practical relevance. An adversary can intercept and store classically encrypted traffic today with the intent of decrypting it once a capable QC becomes available \cite{NIST}. For data with long confidentiality requirements, such as state secrets, medical records, and financial instruments, the effective window of vulnerability is already open. This threat model does not require a quantum computer to exist now; it simply requires that one will exist within the data's confidentiality lifetime.

Two complementary approaches address this threat. PQC replaces vulnerable mathematical assumptions with problems believed to resist quantum attack. NIST finalised its first three PQC standards in August 2024 (FIPS 203, 204, 205), with a fourth algorithm selected in March 2025 \cite{NIST_IR8547}. The accompanying IR-8547 transition roadmap deprecates quantum-vulnerable algorithms by 2030 and disallows them entirely by 2035.

However, as PQC provides \textit{computational} security, its guarantees only hold as long as the underlying mathematical problems remain hard, an assumption that cannot be proven and could be invalidated by future algorithmic breakthroughs. QKD takes a fundamentally different approach. By encoding key material in quantum states and exploiting physical principles such as the no-cloning theorem and measurement disturbance, QKD achieves \textit{information-theoretic} security, remaining secure regardless of an adversary's computational power, now or in the future \cite{Scarani2009, ShorPreskill2000}.

Shannon \cite{Shannon1945} established that information-theoretic secrecy requires a key at least as long as the message, a condition met by the one-time pad but impractical without a secure means of key distribution. QKD is the only known method of generating and distributing such keys at scale without relying on a trusted courier \cite{QKDbible}. Any eavesdropping attempt on the quantum channel inevitably disturbs the transmitted states, producing measurable errors that allow the communicating parties to detect interception and discard compromised key material. This physical guarantee is what distinguishes QKD from all classical and post-quantum alternatives. The role of QKD is therefore not as a competitor to PQC, but as a complementary layer offering a stronger class of security guarantee for applications where long-term confidentiality is critical.

% sources explaining the QM principles for how QKD works
\section{Quantum Mechanical Foundations}

To achieve information-theoretic security, every QKD protocol rests on a small set of quantum mechanical principles. This section introduces the five principles that are exploited by the BB84, B92, and E91 protocols. The treatment of these foundations follows Nielsen and Chuang \cite{QKDbible}, and the pedagogical review by \cite{Scarani2009}. 

\subsection{Qubits \& Superposition}
Consider a classical bit taking one of two definite values, 0 or 1. A Qubit (quantum bit) can exist in a linear combination of both 0 or 1, simultaneously. Using Dirac notation, the general single-qubit state is written as

%
\begin{equation}
\label{eq:qubit}
\ket{\psi} = \alpha\ket{0} + \beta\ket{1},
\end{equation}
%

where $\alpha$ and $\beta$ are complex probability amplitudes that satisfy $|\alpha|^{2}+|\beta|^{2}=1$ \cite{QKDbible}. The two states, $\ket{0}$ and $\ket{1}$ form the \textit{computational basis} (a.k.a. rectilinear or $Z$ basis). Another valid description of a qubit uses the diagonal basis (a.k.a $X$ Basis), defined as

%
\begin{equation}
\label{eq:diag-basis}
\ket{+} = \frac{\ket{0}+\ket{1}}{\sqrt{2}}, \qquad
\ket{-} = \frac{\ket{0}-\ket{1}}{\sqrt{2}}.
\end{equation}
%

In a physical QKD system, qubits are typically encoded in the polarisation of single photons: $\ket{0}$ and $\ket{1}$ may correspond to horizontal and vertical polarisation, while $\ket{+}$ and $\ket{-}$ correspond to $+45^{\circ}$ and $-45^{\circ}$ diagonal polarisation \cite{BB84}. Superposition is what allows a single photon to carry information in more than one basis at once, this property becomes a security asset because any attempt to read the qubit forces a basis choice that irreversibly disturbs the state.

\subsection{Measurement Disturbance}
Measuring a qubit that is in superposition collapses it probabilistically into one of the basis states of the chosen measurement apparatus. If Alice prepares $\ket{+}$ (encoding a logical 1 in BB84's diagonal basis) and Bob measures in the matching $X$ basis, he recovers the correct value with certainty. However, if Bob instead measures in the $Z$ basis, he obtains $\ket{0}$ or $\ket{1}$, each with probability $\frac{1}{2}$, and the original superposition is destroyed \cite{QKDbible}. The $Z$ and $X$ bases are called \textit{conjugate bases} because perfect knowledge in one precludes any information about the other. This is the physical mechanism behind the sifting step in BB84 and B92: only the rounds where Alice and Bob happen to use compatible bases contribute to the key, while mismatched rounds are discarded because Bob's result is uncorrelated with Alice's bit \cite{BB84, B92}.

Critically, measurement disturbance also applies to an eavesdropper (Eve). If Eve intercepts a qubit and measures it in the wrong basis, she collapses the state and re-sends a qubit that no longer matches Alice's preparation. When Alice and Bob later compare a random subset of their sifted key, they detect these errors as an elevated Quantum Bit Error Rate (QBER). An intercept-resend attack on BB84 introduces a characteristic QBER of approximately $25\%$ \cite{Scarani2009, ShorPreskill2000}. % NEW well above the $11\%$ threshold at which no secret key can be


\subsection{No-cloning Theorem}

A natural question is whether Eve could copy the qubit in transit, keep the copy for later analysis, and forward the original undisturbed qubit. The no-cloning theorem, proven independently by Wootters \& Zurek~\cite{WZ1982} and by Dieks \cite{Dieks1982}, rules this attack out. The theorem states that no physical process can produce a perfect copy of an arbitrary unknown quantum state. Formally, there exists no unitary operator $U$ such that $U\ket{\psi}\ket{0} = \ket{\psi}\ket{\psi}$, for all $\ket{\psi}$ \cite{QKDbible}. The proof is compact: suppose such a $U$ did exist for two non-orthogonal states $\ket{\psi}$ and $\ket{\phi}$. Unitarity requires preservation of inner products, giving $\braket{\psi|\phi} = (\braket{\psi|\phi})^{2}$, which is satisfied only when $\braket{\psi|\phi} = 0$ or $1$. Since $\ket{\psi}$ and $\ket{\phi}$ are non-orthogonal by assumption, neither condition holds, producing a contradiction.

No-cloning is the single most important result underpinning QKD. It guarantees that Eve cannot silently duplicate quantum transmissions and is the reason quantum key distribution can offer information-theoretic security rather than the computational security of classical schemes~\cite{ekertQI}. B92 exploits this principle particularly directly, where the protocol encodes bits using only two non-orthogonal states ($\ket{0}$ and $\ket{+}$), no measurement or cloning strategy can distinguish them with certainty \cite{B92}.

\subsection{Heisenberg's Uncertainty}

Closely related to measurement disturbance is Heisenberg's uncertainty principle \cite{Heisenberg1927}. In its information-theoretic form relevant to QKD, the principle states that for two non-commuting observables (such as polarisation measured in
the $Z$ and $X$ bases), simultaneously obtaining precise values for both is impossible. Any measurement that yields information about one observable necessarily introduces uncertainty in the other \cite{QKDbible}. In the context of QKD, this means Eve faces a fundamental trade-off: gaining information about Alice's bit in one basis unavoidably disturbs the conjugate-basis encoding, which manifests as detectable errors. Heisenberg's principle and the no-cloning theorem are logically distinct but operationally complementary; together they establish that any eavesdropping strategy, whether based on measurement, cloning, or a combination of both, will leave a detectable footprint in the quantum channel \cite{Scarani2009}. 

\subsection{Quantum Entanglement \& Bell's Theorem}

Two or more qubits can exist in an \textit{entangled} state that cannot be described as a product of individual qubit states. The canonical example is the Bell state (also called an EPR pair after Einstein, Podolsky and Rosen \cite{EPR1935}): % NEW
%
\begin{equation}
\label{eq:bell-state}
\ket{\Phi^{+}} = \frac{\ket{00}+\ket{11}}{\sqrt{2}}.
\end{equation}
%
If Alice holds the first qubit and Bob the second, a measurement by Alice in the $Z$ basis yielding $\ket{0}$ instantaneously projects Bob's qubit into $\ket{0}$ as well, regardless of the spatial separation between them~\cite{QKDbible}. Entanglement does not permit faster-than-light communication (the individual outcomes are still uniformly random), but it does produce correlations that cannot be explained by any local hidden-variable theory.

Bell \cite{Bell1964} formalised this distinction in 1964 by deriving an inequality that any local realist theory must satisfy. The experimentally testable form due to Clauser, Horne, Shimony and Holt \cite{CHSH1969} (the CHSH inequality) states
%
\begin{equation}
\label{eq:chsh}
|S| = |E(a,b) - E(a,b') + E(a',b) + E(a',b')| \leq 2,
\end{equation}
%
where $E(a,b)$ denotes the correlation of measurement outcomes for settings $a$ and $b$. Quantum mechanics predicts a maximum value of $|S| = 2\sqrt{2} \approx 2.83$, a violation confirmed experimentally by Aspect \cite{Aspect1982}.

Ekert's E91 protocol \cite{E91} directly leverages Bell inequality violation as a built-in eavesdropping test. While BB84 and B92 do not require entanglement, the Shor-Preskill security proof \cite{ShorPreskill2000} establishes a formal equivalence between BB84 and an entanglement-based scheme, meaning the security of prepare-and-measure protocols ultimately derives from the same non-local correlations.

\newpage
% QKD whitepapers
\section{Quantum Key Distribution Protocols}
\begin{figure}[h]
    \centering
    \includegraphics[width=0.6
\linewidth]{images.lit/QKDclassification.png}
    \caption{Classification of QKD Protocol Schemes}
    \label{fig:QKDclassification}
\end{figure}


QKD Protocols can be classified into four basic aspects of use, as depicted in amber in Fig. \ref{fig:QKDclassification}. Each aspect is then further classified based upon the quantum mechanical principles exploited, seen in green. Continuous  Variable (CV) schemes encode data using the continuous variable of a quantum system, such as the amplitude or phase of a light wave. Due to this, CV protocols can be used with common optical equipment and, as such, are less susceptible to noise when deployed practically. In comparison, Discrete Variable (DV) schemes are implemented using discrete quantum states, such as single polarized photons or entangled-pairs of photons \cite{Kumar2025}. Furthermore, QKD protocol structures can classified into three broad paradigms: Prepare \& Measure, Entanglement-Based, or Measurement-Device-Independent schemes. This project focuses on three DV-Protocols: BB84, B92 P\&M protocols, and the Entanglement-Based E91 protocol. 

% edit below this
\subsection{BB84}
\label{sec:bb84}
\begin{figure}[ht]
    \centering
    \includegraphics[width=0.6\linewidth]{images.lit/QuantumAbstract.png}
    \caption{Abstract P\&M QKD Network Model}
    \label{fig:QKDnetworkmodel}
\end{figure}

The BB84 protocol, proposed by Bennett and Brassard in 1984 \cite{BB84}, was the first practical QKD scheme. It is a prepare-and-measure protocol that encodes classical bit values onto single-photon polarisation states drawn from two mutually unbiased bases: the rectilinear basis ($\ket{0},\ket{1}$) and the diagonal basis ($\ket{+},\ket{-}$). The sender, Alice, randomly selects both a bit value and a basis for each qubit, producing one of four possible states:
%
\begin{equation}
  \ket{0},\quad \ket{1},\quad
  \ket{+} = \frac{\ket{0}+\ket{1}}{\sqrt{2}},\quad
  \ket{-} = \frac{\ket{0}-\ket{1}}{\sqrt{2}}.
  \label{eq:bb84-states}
\end{equation}
%


These states are transmitted to Bob over the quantum channel. Bob then independently and randomly chooses a measurement basis for each incoming qubit. Because the two bases are mutually unbiased, a measurement in the wrong basis yields a uniformly random outcome and provides no information about Alice's bit value \cite{QKDbible}.

After transmission is complete, Alice and Bob switch to an authenticated classical channel to perform \textit{sifting}, as seen in Fig. \ref{fig:QKDnetworkmodel}. At this point, Bob announces only his choice of basis for each qubit, and Alice confirms which choices matched her preparation basis. All rounds in which the bases disagree are discarded. In an ideal, noiseless channel, the remaining \textit{sifted key} is perfectly correlated between Alice and Bob, and roughly half the transmitted qubits survive the sifting step \cite{BB84}.

\begin{figure}[h]
    \centering
    \includegraphics[width=0.7\linewidth]{images.lit/bb84operation.png}
    \caption{BB84 Operation Model}
    \label{fig:BB84operation}
\end{figure}

To check for eavesdropping, Alice and Bob sacrifice a random subset of the sifted key and compare bit values over the public channel. Any intercept-resend attack by an eavesdropper, Eve, necessarily disturbs a fraction of the qubits because she cannot simultaneously determine the polarisation in both bases without violating the no-cloning theorem \cite{QKDbible}. A full intercept-resend strategy is expected to produce a Quantum Bit Error Rate (QBER) of approximately 25\%, well above the theoretical security threshold of approximately 11\% \cite{ekertQI}. If the observed QBER exceeds this threshold, the protocol aborts. Otherwise, Alice and Bob proceed through error correction and privacy amplification to yield a final secret key, depicted in Fig. \ref{fig:BB84operation}.


\subsection{B92}
\label{sec:b92}

\begin{figure}[h]
    \centering
    \includegraphics[width=0.6\linewidth]{images.lit/b92basis.png}
    \caption{B92 Measurements Bases}
    \label{fig:B92bases}
\end{figure}

In 1992, Bennett proposed the B92 protocol as a simplification of BB84 \cite{B92}. The central insight is that unconditional security does not require four encoding states; two non-orthogonal states (seen Fig. \ref{fig:B92bases} are sufficient because no measurement can perfectly distinguish between non-orthogonal quantum states \cite{ekertQI}. In B92, Alice encodes bit value 0 as $\ket{0}$ and bit value 1 as $\ket{+}$:

%
\begin{equation}
  \ket{\varphi} =
  \begin{cases}
    \ket{0}  & \text{if } a = 0, \\[4pt]
    \ket{+}  & \text{if } a = 1.
  \end{cases}
  \label{eq:b92-states}
\end{equation}
%

Bob generates his own random bit $a'$ and measures in the $Z$ basis ($\ket{0},\ket{1}$) if $a' = 0$, or in the $X$ basis ($\ket{+},\ket{-}$) if $a' = 1$. When Bob's measurement basis is incompatible with Alice's state, the outcome is conclusive only some of the time, hence, inconclusive results (a.k.a. erasures) are discarded \cite{QKDbible}. Bob publicly announces which rounds yielded a conclusive result, and Alice retains only the corresponding bits. Neither party reveals their basis choice or bit value during this exchange.

Because B92 relies on only two states rather than four, its sifting efficiency is lower than that of BB84, and its tolerance to channel noise is reduced. The theoretical QBER security threshold for B92 is approximately 6.5\%, compared to approximately 11\% for BB84 \cite{ekertQI}. Furthermore, the two-state encoding makes B92 more susceptible to unambiguous state discrimination (USD) attacks, in which Eve exploits the known non-orthogonal structure to occasionally identify Alice's state with certainty while flagging ambiguous cases as losses \cite{B92}. Despite these limitations, B92 remains an important protocol both for its theoretical elegance and because it demonstrates that the minimal resource needed for secure QKD is a pair of non-orthogonal states.

\subsection{E91}
\label{sec:e91}

In contrast to the P\&M paradigm of BB84 and B92, the E91 protocol proposed by Ekert in 1991 \cite{E91} derives its security from quantum entanglement and Bell's theorem. A source prepares pairs of qubits in the singlet Bell state:

%
\begin{equation}
  \ket{\Psi^{-}} = \frac{\ket{01} - \ket{10}}{\sqrt{2}},
  \label{eq:e91-singlet}
\end{equation}
%

and distributes one qubit from each pair to Alice and the other to Bob. Each party independently selects a random measurement orientation from a predetermined set of angles. When Alice and Bob happen to choose the same measurement basis, their outcomes are perfectly anti-correlated, and these matching rounds are used to construct the raw key \cite{E91}.

The rounds in which Alice and Bob chose different bases serve a separate purpose: they are used to evaluate the Clauser-Horne-Shimony-Holt (CHSH) inequality \cite{QKDbible}. For the singlet state, quantum mechanics predicts a CHSH parameter of $S = 2\sqrt{2}$, which maximally violates the classical bound of $S \leq 2$. Any eavesdropping that disturbs the entanglement between the pairs will reduce the observed value of $S$ towards or below the classical limit, thereby revealing Eve's presence \cite{ekertQI, E91}. This connection between eavesdropper detection and Bell inequality violation is the defining feature of E91 and distinguishes it fundamentally from P\&M protocols.

\begin{table}[h]
\centering
\caption{Qualitative Comparison of BB84, B92, and E91}
\label{tab:protocol-comparison}
\footnotesize
\begin{tabularx}{\linewidth}{@{}l X X X@{}}
\toprule
\textbf{Feature} & \textbf{BB84} & \textbf{B92} & \textbf{E91} \\
\midrule
Paradigm & Prepare \& Measure & Prepare \& Measure & Entanglement-Based \\
Encoding states & 4 (two bases) & 2 non-orthogonal & Bell pairs, 3 angles \\
Alice's bases / angles & $\{Z, X\}$ & $\{\ket{0}, \ket{+}\}$ & $\{0, \pi/4, \pi/2\}$ \\
Bob's bases / angles & $\{Z, X\}$ & $\{Z, X\}$ & $\{\pi/4, \pi/2, 3\pi/4\}$ \\
Sifting efficiency & $\sim$50\% & $\sim$25\% & $\sim$22.2\% \\
Threshold ($f_{EC}\!=\!1.0$) & $\sim$11\% & $\sim$8\% & $\sim$11\% \\
Threshold ($f_{EC}\!=\!1.16$) & $\sim$9.8\% & $\sim$6.5\% & $\sim$9.8\% \\
Eve detection & QBER & QBER & CHSH $+$ QBER \\
Full intercept-resend QBER & $\sim$25\% & $\sim$33--38\% & degraded $|S|$ \\
Known vulnerabilities & PNS (WCP sources) & USD attack & Entanglement loss \\
Hardware complexity & Moderate & Lowest & Highest \\
\bottomrule
\end{tabularx}
\end{table}

This entanglement-based approach offers a conceptual advantage as the key bits are not determined until the moment of measurement, so there is no point during transmission at which an encoded classical value exists for Eve to intercept. However, E91 places greater demands on hardware because it requires a source of high-fidelity entangled pairs and the ability to distribute them over the quantum channel while preserving entanglement \cite{ekertQI}. In a simulation context, this translates to the additional complexity of modelling entangled state preparation, Bell-state measurement, and CHSH parameter estimation alongside the standard post-processing pipeline shared with BB84 and B92.



% discuss the basis sifitng & priavcy amplification processes
\section{The Post-Processing Pipeline}
Once the quantum transmission phase of a QKD protocol is complete, the raw measurement outcomes held by Alice and Bob must be transformed into a shared secret key through a sequence of classical post-processing steps. This pipeline is common to all three protocols considered in this project, although the details of sifting differ between prepare-and-measure and entanglement-based schemes.

The first stage is \textit{sifting}, or basis reconciliation. In BB84, Alice and Bob publicly compare their basis choices and discard all rounds where their bases did not match, retaining approximately 50\% of transmitted qubits \cite{BB84}. However, in B92, Bob announces which rounds yielded a conclusive measurement outcome, and only those rounds are kept, producing a sifting efficiency of approximately 25\% \cite{B92}. For E91, sifting separates the rounds into two groups based on whether Alice and Bob selected the same measurement angle. Matching-angle rounds contribute to the raw key, while non-matching rounds are reserved for CHSH parameter estimation \cite{E91}.

Following sifting, Alice and Bob perform \textit{error estimation} by publicly comparing a randomly chosen subset of the sifted key. The fraction of disagreements in this subset provides an estimate of the QBER. If the observed QBER exceeds the protocol-specific security threshold, the key is deemed compromised and the protocol aborts. Otherwise, the sacrificed bits are discarded and the remaining sifted key proceeds to error correction \cite{Scarani2009}.

\textit{Error correction} reconciles the remaining discrepancies between Alice's and Bob's key strings. The CASCADE protocol, introduced by Brassard and Salvail \cite{Cascade1994}, is the most widely used scheme in QKD implementations. CASCADE operates through multiple passes of binary search, where in each pass, the key is divided into blocks of increasing size, and parity checks identify and correct single-bit errors within each block. The multi-pass structure progressively eliminates errors that were missed in earlier rounds. A critical parameter is the error correction efficiency $f_{EC}$, defined as the ratio of information actually leaked during correction to the Shannon-limit minimum. CASCADE achieves $f_{EC} \approx 1.16$ in practice \cite{Cascade1994}, meaning it leaks approximately 16\% more information than the theoretical minimum. This distinction has direct consequences for the achievable secure key rate, as discussed below.

The final stage is \textit{privacy amplification}, which compresses the corrected key to remove any information that may have leaked to an eavesdropper during both the quantum transmission and the error correction phases. Universal hashing functions are applied to map the corrected key of length $n$ to a shorter final key of length $\ell < n$, where $\ell$ is chosen such that Eve's mutual information with the final key is negligibly small \cite{Bennett1995PA}. The amount of compression required depends on both the estimated QBER and the information leaked during error correction.

The GLLP formula \cite{GLLP2004} quantifies how much key material survives this entire pipeline. For BB84 with a single-photon source, the asymptotic secure key rate per sifted bit is given by
%
\begin{equation}
\label{eq:gllp}
R = 1 - h(\text{QBER}) - f_{EC} \cdot h(\text{QBER}),
\end{equation}
%
where $h(\cdot)$ denotes the binary Shannon entropy. Setting $f_{EC} = 1.0$ (the Shannon limit) yields the well-known 11\% QBER threshold for BB84. However, using the realistic CASCADE efficiency of $f_{EC} = 1.16$ shifts this threshold to approximately 9.8\%, a reduction with meaningful engineering consequences for budget calculations. It is worth noting that many simulation studies in the literature default to $f_{EC} = 1.0$, which overestimates the achievable secure key rate. This project uses $f_{EC} = 1.16$ throughout to provide a more accurate representation of practical system performance.

\section{Noise Models for Modelling} 

In any practical QKD deployment, the quantum channel introduces errors that are physically indistinguishable from those caused by eavesdropping. Accurately modelling these channel imperfections is therefore essential for a simulation platform that seeks to evaluate protocol performance under realistic conditions.

The \textit{depolarising channel} is the most commonly used noise model in QKD analysis. It maps an arbitrary single-qubit state $\rho$ to the maximally mixed state with probability $p$, effectively applying a random Pauli error:
%
\begin{equation}
\label{eq:depolarising}
\mathcal{E}_{dep}(\rho) = (1-p)\rho + \frac{p}{3}(X\rho X + Y\rho Y + Z\rho Z).
\end{equation}
%
Because the depolarising channel applies errors isotropically across all axes of the Bloch sphere, it affects all encoding bases equally and is therefore compatible with the assumptions underlying the GLLP key rate formula \cite{GLLP2004}. Physically, it represents a worst-case model of isotropic noise arising from environmental decoherence in the fibre channel.

% block sphere image

Several other noise channels available in Qiskit Aer were considered and excluded from the benchmark. The bit-flip and phase-flip channels act on a single axis of the Bloch sphere and therefore underestimate error on bases orthogonal to that axis; depolarising noise captures the worst-case isotropic behaviour that is appropriate for protocol-level security analysis where the encoding basis should not be treated preferentially. Thermal relaxation requires $T_{1}$/$T_{2}$ time constants that are not physically meaningful for single-photon transmission through fibre. The reset channel models spontaneous emission to the ground state, a process not relevant to photonic QKD systems. This selection prioritises physical relevance to fibre-optic QKD over exhaustive coverage of the Qiskit Aer model library.


\section{QKD Simulation Frameworks}

Several software frameworks exist for simulating quantum communication protocols, each with distinct design philosophies and trade-offs relevant to the goals of this project.

Qiskit Aer \cite{QiskitAer}, the noise simulation backend of IBM's Qiskit framework, provides gate-level circuit simulation with a comprehensive built-in noise model library. Its \texttt{NoiseModel} API allows noise channels to be attached to specific gates rather than applied globally, enabling the channel-marker approach used in this project: an identity gate is inserted between state preparation and measurement to represent the quantum channel, and noise is applied exclusively to this gate. This preserves the trusted-device assumption by ensuring that Alice's preparation and Bob's measurement remain noise-free. No other framework surveyed offered this level of noise insertion control with comparable accessibility. Qiskit Aer's Python-native interface and active maintenance by IBM further supported its selection as the simulation backend for this project.

\section{Research Gaps}
My aim to present the protocols under unified modelling constraints, to obtain a critical analysis of BB84, B92, and E91.

% fix references
The existing QKD simulation literature exhibits several gaps that this project aims to address. The majority of published simulation studies focus on benchmarking a single protocol in isolation. For example, Ma \textit{et al.} \cite{Ma2005} analyse BB84 key rates under various source assumptions, while Tamaki \textit{et al.} \cite{Tamaki2014} examine B92 security bounds independently. Comparative studies that evaluate multiple protocols under identical channel and noise conditions on the same simulation platform are considerably less common, and those that do exist typically compare only two protocols or restrict themselves to a single noise model.

Furthermore, few existing simulation platforms compare prepare-and-measure and entanglement-based protocols side-by-side using a unified benchmarking framework. BB84 and B92 are frequently analysed together due to their shared P\&M structure, but direct comparisons with E91 under matched noise conditions are rare. This limits the ability to draw meaningful conclusions about the relative strengths and weaknesses of these two fundamentally different QKD paradigms.

A further gap concerns the treatment of post-processing parameters in simulation studies. Many published analyses assume Shannon-limit error correction ($f_{EC} = 1.0$), which produces optimistic key rate estimates and overstates the maximum tolerable QBER. In practice, the CASCADE protocol operates at $f_{EC} \approx 1.16$ \cite{Cascade1994}, and this difference has direct consequences for the engineering viability of a given fibre link. Simulations that neglect this factor risk misinforming deployment decisions.

This project addresses these gaps by providing a simulation platform that benchmarks BB84, B92, and E91 under identical, configurable noise conditions with GLLP-based key rate analysis using realistic error correction efficiency. The platform is designed to enable fair comparison between P\&M and entanglement-based QKD protocols on a unified framework.





% ═══════════════════════════════════════════════════════════════════════════════
\chapter{Methodology}
% ═══════════════════════════════════════════════════════════════════════════════

% >>> ADDED SECTION: Simulation Architecture <<<
\section{Simulation Architecture}

% discuss the process of developing the system thru AGILE methodologies, starting priorities lied within correctly implementing protocol logic as opposed to focusing on large system features

\subsection{Modular Package Design}

The simulation platform is implemented as a modular Python package (\texttt{qkd\_sim}), structured across nine source files with clearly separated responsibilities: protocol logic, noise model configuration, eavesdropper simulation, benchmarking orchestration, visualisation, and a command-line interface (CLI). This modular architecture was a deliberate design decision driven by three requirements. First, protocol-agnostic benchmarking: by separating the noise and eavesdropper modules from individual protocol implementations, the same channel conditions can be applied uniformly across BB84, B92, and E91 without code duplication. Second, independent noise model testing: isolating noise configuration into its own module allows channel models to be validated against theoretical predictions before being coupled with protocol logic. Third, reproducibility: decoupling the simulation parameters from the execution code ensures that any experiment can be fully reproduced from its configuration file alone.

The package is organised around an abstract base class, \texttt{QKDProtocol}, which defines the common interface shared by all protocol implementations. Each protocol subclass (BB84, B92, E91) overrides the abstract \texttt{run()} method while inheriting the shared post-processing logic. A companion dataclass, \texttt{QKDResult}, encapsulates the outputs of every protocol run, including Alice's and Bob's sifted keys, the computed QBER, sifting efficiency, and the secure key rate derived from the GLLP formula. This object-oriented design ensures that the benchmarking module can treat all protocols interchangeably, enabling fair comparison under identical conditions.

% NOTE: Replace with your actual source listing or figure
\begin{figure}[ht]
    \centering
    \fbox{\parbox{0.9\linewidth}{\centering\textit{[Placeholder: Package architecture diagram showing module dependencies: config $\rightarrow$ protocols (bb84, b92, e91) $\leftarrow$ base; noise $\rightarrow$ backend; eve $\rightarrow$ protocols; benchmark $\rightarrow$ protocols + noise; plotter $\leftarrow$ benchmark; \_\_main\_\_ $\rightarrow$ config + benchmark + plotter]}}}
    \caption{Module dependency structure of the \texttt{qkd\_sim} package. Arrows indicate import dependencies.}
    \label{fig:architecture}
\end{figure}

\subsection{Configuration System}
Initially, simulation parameters were hardcoded into the system to accelerate the process of implementing the protocol logic. Through AGILE development methodologies, I was able to expand the capability of the system by implementing configurable YAML files to load simulation parameters into the system. YAML was selected over JSON for its support of inline comments, which allows each parameter to be annotated with its physical meaning and valid range, and over TOML for its more natural handling of nested structures. The configuration is loaded into a \texttt{SimConfig} dataclass that validates the execution mode against the four supported modes: \texttt{single} (one protocol trial with fixed parameters), \texttt{sweep} (noise strength sweep), \texttt{eve\_sweep} (eavesdropper interception rate sweep), and \texttt{protocol\_comparison} (identical noise sweep applied to multiple protocols simultaneously). A representative configuration snippet is shown in \cref{lst:yaml-config}.

% NOTE: Replace with \lstinputlisting pointing to an actual config file
\begin{lstlisting}[language=Python, caption={Example YAML configuration for a BB84 noise sweep}, label={lst:yaml-config}]
# Placeholder: insert actual YAML config snippet
# mode: sweep
# protocol: bb84
# n_qubits: 1000
# n_trials: 50
# noise_type: depolarizing
# noise_min: 0.0
# noise_max: 0.30
# noise_steps: 21
# f_ec: 1.16
# e91_channel_topology: both
\end{lstlisting}

This configuration-driven approach directly supports objective O5.3 (reproducibility) by ensuring that every simulation run is fully parameterised and traceable.

\subsection{Execution Pipeline}

The simulation follows a linear execution pipeline, illustrated in \cref{fig:pipeline}: a YAML configuration file is parsed by the \texttt{SimConfig} loader to select the protocol, channel parameters, and execution mode; the CLI dispatcher in \texttt{\_\_main\_\_.py} delegates to the appropriate runner function; Alice prepares quantum states according to the selected protocol; an identity gate is inserted as a channel marker to which noise models are attached; the optional eavesdropper module intercepts a configurable fraction of qubits; Bob performs his measurement; classical sifting extracts the raw key; and the post-processing logic computes the QBER, applies the GLLP key rate formula, and returns a \texttt{QKDResult} object. For sweep and comparison modes, the \texttt{BenchmarkRunner} orchestrates repeated trials across parameter ranges and aggregates the results into \texttt{BenchmarkData} objects containing mean and standard deviation arrays, which are then passed to the \texttt{QKDPlotter} for visualisation.

\begin{figure}[ht]
    \centering
    \fbox{\parbox{0.9\linewidth}{\centering\textit{[Placeholder: Execution pipeline block diagram: YAML Config $\rightarrow$ SimConfig $\rightarrow$ CLI Dispatcher $\rightarrow$ Protocol Selection $\rightarrow$ State Preparation $\rightarrow$ Channel (id gate + noise) $\rightarrow$ Eve (optional) $\rightarrow$ Measurement $\rightarrow$ Sifting $\rightarrow$ Post-Processing $\rightarrow$ QKDResult $\rightarrow$ BenchmarkData $\rightarrow$ Visualisation]}}}
    \caption{Simulation execution pipeline. Noise is applied exclusively to the channel stage via the identity gate marker, preserving the trusted-device assumption.}
    \label{fig:pipeline}
\end{figure}
% >>> END ADDED SECTION <<<


% >>> ADDED SECTION: Protocol Implementations <<<
\section{Protocol Implementations}
\label{sec:protocol-implementations}

This section describes the implementation of each QKD protocol within the simulation platform. All three protocols share the same post-processing pipeline and noise application mechanism; the differences lie in state preparation, measurement, and sifting logic.

\subsection{BB84 Protocol}

The BB84 implementation follows the four-state, two-basis scheme described in \cref{sec:bb84}. For each of $n$ qubits, Alice randomly selects a bit value $a_i \in \{0,1\}$ and a basis $\hat{a}_i \in \{Z, X\}$ using NumPy's pseudorandom number generator, then prepares the corresponding quantum state on a single-qubit Qiskit circuit. If the selected bit is 1, a Pauli-X gate is applied to flip the qubit from $\ket{0}$ to $\ket{1}$. If the diagonal basis is selected, a Hadamard gate is subsequently applied to rotate into the $X$ basis. An identity gate (\texttt{qc.id(0)}) is appended after state preparation to serve as the channel marker, to which noise models are attached during execution. Bob independently selects a random measurement basis $\hat{b}_i$; if he chooses the diagonal basis, a Hadamard gate is applied before the computational-basis measurement. All $n$ circuits are executed in a single batch call to the Qiskit Aer backend with \texttt{shots=1} and \texttt{memory=True} to retrieve individual outcomes.

The sifting procedure retains only the rounds where $\hat{a}_i = \hat{b}_i$, implemented as a Boolean mask comparison of the two basis arrays. On an ideal channel, this produces a sifting efficiency of approximately 50\% and a QBER of 0\%, serving as the primary validation targets for this protocol. The key rate is computed as the ratio of sifted key length to total transmitted qubits.

% NOTE: Replace firstline/lastline with actual line numbers from bb84.py
\lstinputlisting[
    language=Python,
    caption={BB84 state preparation and channel-marker insertion (\texttt{bb84.py}).},
    label={lst:bb84-prep},
    firstline=1,
    lastline=1
]{../code/qkd_sim/protocols/bb84.py}  % TODO: set correct path + line range

The high-level flow of the BB84 key sifting procedure is presented in \cref{alg:bb84}.

\begin{algorithm}[H]
    \caption{BB84 Key Sifting Procedure}\label{alg:bb84}
    \KwIn{Alice's bits $a_i$, bases $\hat{a}_i$; Bob's measurements $b_i$, bases $\hat{b}_i$}
    \KwOut{Sifted key $K$}
    $K \leftarrow \emptyset$\;
    \For{$i \leftarrow 1$ \KwTo $n$}{
        Alice publishes $\hat{a}_i$ over classical channel\;
        Bob publishes $\hat{b}_i$ over classical channel\;
        \If{$\hat{a}_i = \hat{b}_i$}{
            Append $b_i$ to $K$ \Comment{Matching basis -- keep bit}
        } 
    }
    \Return{$K$}
\end{algorithm}


\subsection{B92 Protocol}

The B92 implementation uses only two non-orthogonal encoding states: $\ket{0}$ for bit value 0 and $\ket{+}$ for bit value 1. Alice prepares $\ket{0}$ directly (no gate applied) or applies a Hadamard gate to produce $\ket{+}$, followed by the identity gate channel marker (\texttt{qc.id(0)}).

Bob's measurement logic is the most distinctive aspect of the B92 implementation. Bob generates a random bit $b_i \in \{0, 1\}$ and measures in the $Z$ basis if $b_i = 0$ or the $X$ basis if $b_i = 1$, applying a Hadamard gate before measurement in the latter case. An outcome is \textit{conclusive} only if, given the basis chosen, the observed result is inconsistent with one of Alice's two possible preparations. Measuring $\ket{0}$ in the $Z$ basis can only yield $0$ and is therefore ambiguous, as is measuring $\ket{+}$ in the $X$ basis (which also yields $0$ under the Hadamard-then-$Z$ implementation). The conclusive events are the \emph{complementary} pairings: a $Z$-basis outcome of $1$ rules out a $\ket{0}$ preparation and therefore identifies Alice's bit as $1$; an $X$-basis outcome of $1$ (equivalent to projecting onto $\ket{-}$) rules out $\ket{+}$ and identifies Alice's bit as $0$. Inconclusive rounds (outcome $0$) are discarded.

The sifting step in B92 therefore differs from BB84: Bob announces which rounds were conclusive, not which basis he used. Neither party reveals their bit value or basis choice during this exchange. Bob's sifted key is inferred as $1 - b_i$ (the complement of Bob's measurement basis for the conclusive rounds), while Alice retains her original bit values for those rounds. This sifting mechanism produces a theoretical sifting efficiency of approximately 25\%, half that of BB84. The reduced efficiency, combined with the two-state encoding, results in a lower noise tolerance.

% NOTE: Replace firstline/lastline with actual line numbers from b92.py
\lstinputlisting[
    language=Python,
    caption={B92 conclusive-measurement decoding and sifting (\texttt{b92.py}).},
    label={lst:b92-sift},
    firstline=1,
    lastline=1
]{../code/qkd_sim/protocols/b92.py}  % TODO: set correct path + line range


\subsection{E91 Protocol}

The E91 implementation models the entanglement-based protocol described in \cref{sec:e91}. For each round, a two-qubit Qiskit circuit is initialised and a singlet Bell state $\ket{\Psi^{-}} = (\ket{01} - \ket{10})/\sqrt{2}$ is prepared by applying the following gate sequence: a Hadamard gate on qubit 0, a CNOT gate with qubit 0 as control and qubit 1 as target, a Pauli-X gate on qubit 0, and a Pauli-Z gate on qubit 0. This sequence produces the singlet state rather than the more commonly used $\ket{\Phi^+}$, matching the original E91 specification \cite{E91}. Qubit 0 is designated as Alice's and qubit 1 as Bob's.

The implementation supports two channel topologies via a configurable parameter. In the \texttt{bob} topology, an identity gate channel marker is applied only to Bob's qubit (\texttt{qc.id(1)}), representing a scenario where noise acts on the distribution channel to Bob while Alice's qubit remains local. In the \texttt{both} topology, identity gates are applied to both qubits (\texttt{qc.id(0)} and \texttt{qc.id(1)}), modelling noise on both distribution paths. This distinction affects the relationship between noise strength and observed QBER: under depolarising noise with parameter $p$, the \texttt{bob} topology produces $\text{QBER} = p/2$ (matching BB84), while the \texttt{both} topology produces $\text{QBER} = p - p^2/2$.

Alice and Bob each independently select a measurement angle from predetermined sets. Alice's angles are drawn from $\{0, \pi/4, \pi/2\}$ and Bob's from $\{\pi/4, \pi/2, 3\pi/4\}$, each chosen uniformly at random per round. Measurement in an arbitrary angle $\theta$ is implemented by applying a rotation gate $R_y(-\theta)$ before a $Z$-basis measurement. A predefined $3 \times 3$ sift table, indexed by Alice's and Bob's angle choices, classifies each round as contributing to the key, to the CHSH parameter estimation, or as discarded. Of the nine possible angle-pair combinations, two contribute to the key (matching angles that yield anti-correlated outcomes), four contribute to CHSH estimation, and three are discarded, producing a theoretical sifting efficiency of $2/9 \approx 22.2\%$.

The CHSH parameter is computed as
%
\begin{equation}
\label{eq:chsh-computation}
S = +E(a_1, b_1) - E(a_1, b_3) + E(a_3, b_1) + E(a_3, b_3),
\end{equation}
%
where $E(a_i, b_j)$ is the correlation coefficient for the angle pair $(a_i, b_j)$, calculated as the expectation value of the product of Alice's and Bob's outcomes (mapped to $\pm 1$ via $1 - 2 \times \text{outcome}$) across all rounds using that pair. The implementation is validated by confirming that $|S| \rightarrow 2\sqrt{2} \approx 2.828$ in the absence of noise or eavesdropping, consistent with the maximal quantum violation of the CHSH inequality. Under depolarising noise, the theoretical CHSH value is $2\sqrt{2}(1-p)$ for the \texttt{bob} topology and $2\sqrt{2}(1-p)^2$ for the \texttt{both} topology.

% NOTE: Replace firstline/lastline with actual line numbers from e91.py
\lstinputlisting[
    language=Python,
    caption={E91 Bell-state preparation and CHSH computation (\texttt{e91.py}).},
    label={lst:e91-chsh},
    firstline=1,
    lastline=1
]{../code/qkd_sim/protocols/e91.py}  % TODO: set correct path + line range
% >>> END ADDED SECTION <<<


% >>> ADDED SECTION: Modelling Assumptions <<<
\section{Modelling Assumptions}
\label{sec:modelling-assumptions}

The simulation platform operates under a set of assumptions that define the scope and limitations of the model. Stating these explicitly is important because they determine which physical effects are captured and which are deliberately excluded.

The platform assumes a perfect single-photon source where each transmission round involves exactly one photon, with no multi-photon pulses or vacuum events. Real-world QKD implementations typically use weak coherent pulse sources, which emit multi-photon pulses with non-zero probability. Multi-photon events enable photon-number-splitting attacks that are not modelled here.

All noise is assumed to originate in the quantum channel. Alice's state preparation and Bob's measurement apparatus are treated as perfect systems. This is enforced in the implementation by attaching noise models exclusively to the identity gate channel marker, ensuring that preparation and measurement operations remain noise-free. This assumption is standard in the security proofs of BB84, B92, and E91 \cite{ShorPreskill2000, Scarani2009}.

The authenticated classical channel used for sifting, error estimation, and error correction is assumed to be tamper-proof and perfectly secure. While this is a standard assumption in QKD security analysis, it is worth noting that it requires a pre-shared authentication key, creating a problem that is outside the scope of this simulation.

Finally, the model does not account for detector efficiency mismatch, timing jitter, or other side-channel imperfections. These effects are relevant to device-security analysis but are excluded here to maintain focus on the protocol-level performance as the primary objective.


% >>> ADDED SECTION: Noise Models and Channel Simulation <<<
\section{Noise Models and Channel Simulation}
\label{sec:noise-impl}

This section describes how the noise models introduced in the literature review are implemented within the simulation platform.

\subsection{Channel Marker Approach}

Noise is applied to the quantum channel using Qiskit Aer's \texttt{NoiseModel} API. The \texttt{create\_backend()} function in the noise module constructs an \texttt{AerSimulator} instance with a \texttt{NoiseModel} that attaches the specified quantum error exclusively to the \texttt{'id'} gate. An identity gate (\texttt{qc.id()}) is inserted into each single-qubit circuit between Alice's state preparation and Bob's measurement. This gate performs no logical operation but serves as a marker to which noise channels are attached. When the circuit is executed on the Qiskit Aer simulator with the configured \texttt{NoiseModel}, noise is applied exclusively to this identity gate.

This approach is critical for maintaining the trusted-device assumption. If noise were applied globally to all gates, the Hadamard gates used in Alice's preparation and Bob's measurement would also be affected, conflating device imperfection with channel noise and artificially inflating the observed QBER. By isolating noise application to the channel marker, the simulation accurately represents a scenario where only the quantum transmission link is imperfect.

The noise module supports the ideal \& depolarising channel for the benchmarking analysis, as justified in the literature review.

\subsection{Eavesdropper Model}

The eavesdropper simulation is implemented in the \texttt{EveInterceptor} class, which performs an intercept-resend attack with a configurable interception probability $p \in [0, 1]$. For each qubit in transit, Eve intercepts it with probability $p$, selects a random measurement basis, and measures the intercepted qubit by appending the appropriate gates and executing the circuit on the simulator backend. Eve then prepares a replacement qubit in the state corresponding to her measurement outcome and basis choice, including an identity gate channel marker for any noise on the Eve-to-Bob link, before forwarding it to Bob. All intercepted qubits are measured in a single batch call to the simulator for computational efficiency.

For BB84 and B92, Eve's random basis choice means she selects the correct basis with probability $\frac{1}{2}$. When she guesses incorrectly, her measurement collapses the qubit into a state incompatible with Alice's preparation, introducing a 50\% error probability on those rounds. The net effect for BB84 is a QBER contribution of $p \times 0.5 \times 0.5 = 0.25p$, reaching the characteristic 25\% at full interception ($p = 1$). For E91, Eve's interception degrades the entanglement fidelity of the Bell pairs, reducing the observed CHSH parameter $|S|$ towards the classical bound of 2. This provides an independent, physics-based eavesdropping detection mechanism that does not rely on QBER estimation.

The interception probability is swept from 0 to 1 in benchmarking experiments to generate the QBER-versus-interception-rate curves that form a central part of the results analysis.
% >>> END ADDED SECTION <<<


% >>> ADDED SECTION: Post-Processing Implementation <<<
\section{Post-Processing Implementation}

The post-processing pipeline is shared across all three protocols and consists of error estimation, secure key rate calculation, and mutual information computation.

\subsection{Error Estimation}

After sifting, the QBER is computed directly as the ratio of bit disagreements between Alice's and Bob's sifted keys to the total sifted key length. This computation is performed within the \texttt{QKDResult} dataclass, which stores both sifted keys and exposes the QBER as a derived property. The implementation does not explicitly sacrifice a subset of the sifted key for QBER estimation, as the simulation has access to both keys in their entirety; in a physical deployment, a configurable fraction would be publicly compared and discarded.

\subsection{GLLP Key Rate Calculation}

The secure key rate is computed using the GLLP formula presented in \cref{eq:gllp}. The binary Shannon entropy function is implemented within the \texttt{QKDResult} dataclass as a private method:
%
\begin{equation}
\label{eq:binary-entropy}
h(x) = -x \log_2(x) - (1 - x) \log_2(1 - x),
\end{equation}
%
with the convention that $0 \log_2(0) = 0$. The secure key rate property computes the secret fraction as $\max(0,\; 1 - h(\text{QBER}) - f_{EC} \cdot h(\text{QBER}))$ and multiplies it by the sifting efficiency (key rate) to yield the net secure key rate per transmitted qubit. The error correction efficiency is set to $f_{EC} = 1.16$ throughout, corresponding to a practical CASCADE implementation. This choice shifts the BB84 QBER security threshold from the idealised $\sim$11\% (at $f_{EC} = 1.0$) to $\sim$9.8\%, and reduces the B92 threshold to $\sim$6.5\%. These thresholds are the QBER values at which $R = 0$, beyond which no secret key can be extracted.

\subsection{Mutual Information}

For the mutual information analysis, the platform computes the Alice-Bob mutual information $I(A{:}B) = 1 - h(\text{QBER})$ and the Alice-Eve mutual information $I(A{:}E) = h(\text{QBER})$ as functions of the observed QBER. The crossover point where $I(A{:}E) \geq I(A{:}B)$ corresponds to the threshold beyond which secure key generation is impossible. This computation provides an alternative visualisation of the security threshold that complements the GLLP key rate analysis.
% >>> END ADDED SECTION <<<


% >>> ADDED SECTION: Performance Metrics <<<
\section{Performance Metrics}
\label{sec:methodology-metrics}

The benchmarking framework evaluates protocol performance using the following metrics, each of which is computed per simulation run and aggregated across multiple runs to produce confidence intervals.

The \textit{Quantum Bit Error Rate} (QBER) is defined as the ratio of bit errors in the sifted key to the total sifted key length. It is the primary indicator of channel quality and eavesdropper presence.

The \textit{sifting efficiency} is the ratio of sifted key length to total transmitted qubits. Theoretical targets are 50\% for BB84, 25\% for B92, and $2/9 \approx 22.2\%$ for E91 based on the sift table configuration.

The \textit{secure key rate} is the output of the GLLP formula (\cref{eq:gllp}), expressed as bits of secure key per sifted bit. A positive key rate indicates that secure key generation is feasible at the given noise level.

For E91, the \textit{CHSH parameter} $S$ is computed from the measurement correlations across non-matching-basis rounds, as defined in \cref{eq:chsh-computation}. A value of $|S| > 2$ confirms the presence of quantum correlations; degradation towards $|S| = 2$ indicates potential eavesdropping.

Each protocol class also provides a \texttt{theoretical\_qber()} method that returns the analytical QBER as a function of noise type and strength. For BB84, the theoretical QBER under depolarising noise is $p/2$. For B92, the depolarising theoretical QBER is $p/(1+p)$. These analytical predictions are overlaid on simulation data in the benchmarking plots as a validation mechanism.

\subsection{Statistical Methodology}

Each data point in the benchmarking results represents the mean of $N$ independent simulation runs (default: 30, configurable via the YAML configuration), each using $n$ qubits per run (default: 1000, configurable). The \texttt{BenchmarkRunner} collects the QBER, key rate, and secure key rate from each trial and computes the mean and standard deviation across trials. Error bars in the visualisations represent the standard error of the mean (SEM), computed as $\text{SEM} = \sigma / \sqrt{N}$, providing a measure of the uncertainty in the mean estimate rather than the spread of individual trials.

% NOTE: Update N and n values above to match your actual simulation parameters.
% >>> END ADDED SECTION <<<


% >>> ADDED SECTION: Results & Discussion Chapter <<<
% ═══════════════════════════════════════════════════════════════════════════════
\chapter{Results \& Discussion}
% ═══════════════════════════════════════════════════════════════════════════════

This chapter presents the simulation results, structured to follow the project objectives. The analysis begins with ideal-channel validation of the three implementations (O5.1), progresses through eavesdropper detection in prepare-and-measure protocols (O3.2) and in the entanglement-based protocol (O2.3), compares noise tolerance across all three schemes (O3.1, O5.3), and closes with an information-theoretic capacity benchmark (O4.2) and a discussion of deployment feasibility (O6.1). Each observation is tied back to the analytical predictions established in the literature review so that the data both verify the simulator and serve as a quantitative basis for protocol comparison.

Unless stated otherwise, every data point presented in this chapter is the mean of $N = 20$ independent trials of $n = 1000$ qubits each (O5.3), executed on the Qiskit Aer QASM simulator with the noise model attached to the channel-marker identity gate as described in \cref{sec:noise-impl}. Error bars show the standard error of the mean, $\sigma/\sqrt{N}$.

% ─────────────────────────────────────────────────────────────────────────────
\section{Ideal Channel Validation}
\label{sec:results-ideal}

Before interpreting any noise-channel or adversarial result, the simulator must be shown to reproduce the analytical predictions of each protocol on an ideal channel. \Cref{tab:validation} collects the four pass/fail metrics defined in objective O5.1. The analytically expected values are listed alongside the simulated means; every simulated value is expected to lie within $\pm 2\%$ of its theoretical target, satisfying the validation criteria set in \cref{sec:methodology-metrics}.

\begin{table}[ht]
    \centering
    \caption{Ideal Channel Validation Results ($n=1000$, $N=20$)}
    \label{tab:validation}
    \footnotesize
    \begin{tabularx}{\linewidth}{@{}l c c X@{}}
        \toprule
        \textbf{Metric} & \textbf{Target} & \textbf{Simulated} & \textbf{Objective} \\
        \midrule
        BB84 sifting efficiency (\%)     & $50.0$  & \textit{[insert]} & O5.1 \\
        B92 sifting efficiency  (\%)     & $25.0$  & \textit{[insert]} & O5.1 \\
        E91 sifting efficiency  (\%)     & $22.2$  & \textit{[insert]} & O5.1 \\
        BB84 ideal-channel QBER (\%)     & $0.0$   & \textit{[insert]} & O5.1 \\
        B92 ideal-channel QBER  (\%)     & $0.0$   & \textit{[insert]} & O5.1 \\
        E91 ideal-channel QBER  (\%)     & $0.0$   & \textit{[insert]} & O5.1 \\
        E91 CHSH parameter $|S|$         & $2\sqrt{2}\!\approx\!2.828$ & \textit{[insert]} & O2.3 \\
        \bottomrule
    \end{tabularx}
\end{table}

The ideal-channel sifting efficiencies confirm the protocol-level accounting developed in \cref{sec:protocol-implementations}. BB84 retains half of all transmitted qubits because Alice and Bob each choose one of two bases uniformly and independently; B92 retains a quarter because Bob's conclusive-measurement criterion fires on half of the basis-matched rounds; and E91 retains two of the nine possible Alice--Bob angle pairings, giving $2/9\!\approx\!22.2\%$. The zero QBER across all three protocols verifies that the identity channel-marker and the post-processing pipeline introduce no spurious bit errors. For E91 the CHSH parameter converges to the Tsirelson bound of $2\sqrt{2}$, confirming both the fidelity of the singlet preparation circuit and the correctness of the correlation accumulator used to estimate $|S|$. Together these checks give a baseline of confidence that every subsequent deviation from ideal behaviour is attributable either to the injected noise model or to the simulated eavesdropper, and not to an implementation artefact.

% ─────────────────────────────────────────────────────────────────────────────
\section{Eavesdropper Detection in Prepare-and-Measure Protocols}
\label{sec:results-eve-pm}

\subsection{B92 under the Intercept-Resend Attack}

\Cref{fig:b92-eve} shows the QBER of B92 as a function of the eavesdropper interception probability $p \in [0, 1]$ on an otherwise noiseless channel, addressing O3.2.

\begin{figure}[h]
    \centering
    \includegraphics[width=\linewidth]{images.sim/b92_qber_eve.png}
    \caption{B92 QBER versus Eve interception rate ($n=1000$, $N=20$). The horizontal dashed line marks the $f_{EC}\!=\!1.16$ security threshold of 6.5\%; the vertical dotted line marks the interception rate at which the threshold is first exceeded, $p\!\approx\!0.16$.}
    \label{fig:b92-eve}
\end{figure}

The curve is monotonic and approximately linear across the sweep, rising from $\sim\!0\%$ at $p=0$ to $\sim\!37\%$ at $p=1$. The security threshold at 6.5\% is crossed at $p\!\approx\!0.16$, meaning that B92 must abort if even 16\% of qubits have been intercepted. The $p=1$ endpoint of $\sim\!37\%$ agrees closely with the analytical prediction for this attack strategy. Because Eve mimics Bob's measurement logic (random $Z/X$ basis, then conclusive-outcome inference via the same $1 - \text{basis}$ rule used in \texttt{B92Protocol.\_post\_process()}), she obtains a conclusive outcome with probability $1/4$ and is forced to guess otherwise, giving an overall probability of $5/8$ of correctly identifying Alice's bit and $3/8$ of re-encoding the wrong B92 state. When Eve's replay is wrong, Bob's conclusive outcome identifies Eve's bit rather than Alice's, so the QBER on the sifted subset converges to $3/8\!=\!37.5\%$ at $p\!=\!1$, independent of Bob's basis choice by symmetry. The curve is therefore the straight line $\text{QBER}(p)\!\approx\!3p/8$, and the 6.5\% threshold at $p\!\approx\!6.5/37.5\!\approx\!0.173$ matches the observed crossing. The factor-of-1.5 increase in slope relative to BB84's $p/4$ is a quantitative expression of why non-orthogonal two-state encoding is more fragile against intercept-resend than four-state conjugate-basis encoding.

\subsection{BB84 under the Intercept-Resend Attack}

BB84's response to intercept-resend is analysed in two stages: first on a pristine quantum channel to validate the theoretical prediction, then with a realistic channel-noise floor.

\Cref{fig:bb84-eve-pristine} shows the pristine result: QBER versus Eve's interception probability with no depolarising noise applied to the channel.

\begin{figure}[h]
    \centering
    \includegraphics[width=\linewidth]{images.sim/bb84_ideal_eve_qber.png}
    \caption{BB84 QBER versus Eve interception rate on a pristine channel ($n=1000$, $N=20$). The horizontal dashed line marks the Shor--Preskill 11\% bound; the vertical dotted line marks the crossing at $p\!\approx\!0.44$.}
    \label{fig:bb84-eve-pristine}
\end{figure}

The curve rises monotonically and approximately linearly from $\sim\!0\%$ at $p=0$ to $\sim\!26\%$ at $p=1$, in close agreement with the analytical prediction $\text{QBER}(p) = p/4$ derived from the intercept-resend geometry: Eve guesses the basis correctly with probability $1/2$, and when she guesses incorrectly her replay collapses Bob's measurement into a state that disagrees with Alice's bit with probability $1/2$, yielding a net error contribution of $p \cdot \tfrac{1}{2} \cdot \tfrac{1}{2} = p/4$~\cite{ShorPreskill2000, Scarani2009}. The observed $p=1$ endpoint of $\sim\!26\%$ is within statistical error of the theoretical ceiling of $25\%$. The 11\% Shor--Preskill bound is crossed at $p\!\approx\!0.44$; the tighter $f_{EC}\!=\!1.16$ practical threshold of 9.8\% is crossed earlier at $p\!\approx\!0.39$. BB84 therefore tolerates Eve intercepting up to $\sim\!39\%$ of qubits before the protocol must abort on a pristine link.

\Cref{fig:bb84-eve-noise} repeats the sweep with a $5\%$ depolarising noise floor applied to the channel. This is the operationally meaningful scenario: a deployed BB84 link will always carry some residual channel noise, and Eve must be detected against that background.

\begin{figure}[h]
    \centering
    \includegraphics[width=\linewidth]{images.sim/bb84_qber_noise_eve.png}
    \caption{BB84 QBER versus Eve interception rate with a $p_{\mathrm{dep}} = 0.05$ depolarising noise floor ($n=1000$, $N=20$). The horizontal dashed line is the Shor--Preskill 11\% bound; the vertical dotted line marks the observed crossing at $p\!\approx\!0.34$. The practical $f_{EC}\!=\!1.16$ threshold at 9.8\% is crossed earlier still at $p\!\approx\!0.29$.}
    \label{fig:bb84-eve-noise}
\end{figure}

Three features of this result merit comment. First, the $p=0$ intercept gives a baseline QBER of $\sim\!2.4\%$, in excellent agreement with the analytical prediction $\text{QBER}_{\mathrm{dep}} = p_{\mathrm{dep}}/2 = 2.5\%$ for BB84 under a depolarising channel (\cref{sec:methodology-metrics}). Second, the full-interception endpoint at $p=1$ produces $\text{QBER}\!\approx\! 28\%$, which decomposes cleanly as $\sim\!25\%$ from the intercept-resend attack plus the $\sim\!2.5\%$ floor from the residual depolarising noise, confirming that the channel-marker abstraction correctly composes independent error processes. Third, the 11\% Shor--Preskill bound is breached at $p\!\approx\!0.34$, ten percentage points of interception earlier than on the pristine channel, and the practical 9.8\% threshold at $p\!\approx\!0.29$. The contrast between \cref{fig:bb84-eve-pristine} and \cref{fig:bb84-eve-noise} is instructive: moving from an idealised study ($f_{EC}\!=\!1.0$, zero noise, threshold crossing at $p\!\approx\!0.44$) to a realistic one ($f_{EC}\!=\!1.16$, 5\% floor, threshold crossing at $p\!\approx\!0.29$) lowers the tolerable-interception margin by roughly a third --- a degradation that is entirely obscured when simulation studies default to idealised assumptions.

\subsection{Comparative Discussion}

Comparing the two P\&M protocols on matched pristine channels yields a concrete engineering statement: BB84 tolerates roughly \emph{two and a half times} the interception rate of B92 before the protocol must abort ($p\!\approx\!0.39$ versus $p\!\approx\!0.16$ at $f_{EC}\!=\!1.16$). This ratio is not an accident of numerics; it reflects the combined effect of B92's smaller alphabet (which amplifies the error signature of any basis mismatch) and its lower security threshold (which leaves less error budget before abort). The result quantifies the security cost of B92's encoding simplicity in operationally meaningful terms.

\Cref{tab:eve-ranking} collects the threshold-crossing results from the three sub-experiments above.

\begin{table}[h]
    \centering
    \caption{P\&M Protocol Eavesdropper-Detection Thresholds}
    \label{tab:eve-ranking}
    \footnotesize
    \begin{tabularx}{\linewidth}{@{}l c c c c@{}}
        \toprule
        \textbf{Scenario} & \textbf{SP-bound $p$} & \textbf{$f_{EC}\!=\!1.16$ $p$} & \textbf{QBER at $p\!=\!1$} & \textbf{Figure} \\
        \midrule
        BB84, pristine channel   & $0.44$ & $0.39$ & $\sim\!26\%$ & \ref{fig:bb84-eve-pristine} \\
        BB84, $5\%$ noise floor  & $0.34$ & $0.29$ & $\sim\!28\%$ & \ref{fig:bb84-eve-noise} \\
        B92, pristine channel    & ---    & $0.16$ & $\sim\!37\%$ & \ref{fig:b92-eve} \\
        \bottomrule
    \end{tabularx}
\end{table}

% ─────────────────────────────────────────────────────────────────────────────
\section{Entanglement-Based Detection: CHSH Degradation}
\label{sec:results-e91}

E91 detects adversarial activity through a physical signature rather than an error statistic: any loss of entanglement fidelity between the distributed Bell pairs reduces the observed CHSH correlator $|S|$ toward the classical bound of 2. \Cref{fig:e91-chsh} shows $|S|$ and the companion QBER as functions of depolarising noise strength in the \texttt{both} channel topology, in which each qubit of the singlet traverses an independent noisy channel.

\begin{figure}[h]
    \centering
    \includegraphics[width=\linewidth]{images.sim/e91_chsh_qber.png}
    \caption{E91 CHSH parameter $|S|$ (green, left axis) and QBER (red, right axis) versus depolarising noise strength $p$ under the \texttt{both} channel topology ($n=1000$ photon pairs, $N=20$). The horizontal dashed lines mark the classical bound $|S|\!=\!2$ and the Tsirelson bound $|S|\!=\!2\sqrt{2}$.}
    \label{fig:e91-chsh}
\end{figure}

At $p=0$ the measured $|S|\!\approx\!2.82$ saturates the Tsirelson bound, confirming that the simulator reproduces the maximal quantum violation of the CHSH inequality for the singlet state~\cite{CHSH1969, QKDbible}. As noise is introduced, $|S|$ degrades along a curve that is well described by the analytical prediction for symmetric depolarising noise on both sides of the channel,
%
\begin{equation}
\label{eq:chsh-both}
|S|(p) = 2\sqrt{2}\,(1-p)^{2},
\end{equation}
%
which gives $|S|(0.10) = 2.29$, $|S|(0.15) = 2.04$, and $|S|(0.20) = 1.81$. Each of these analytical values lies within the simulated error bars. The classical bound $|S|=2$ is crossed at $p\!\approx\!0.17$, the point at which E91's Bell-test-based eavesdropping alarm would fire.

The QBER curve grows in parallel according to the \texttt{both}-topology prediction $\text{QBER}(p) = p - p^{2}/2$, implemented directly in \texttt{E91Protocol.theoretical\_qber()} and derived from the loss of anti-correlation when both halves of the singlet are independently depolarised before measurement. The prediction passes through $9.5\%$ at $p=0.10$ and $18.0\%$ at $p=0.20$. Significantly, the CHSH alarm ($|S|\le 2$) and the QBER alarm ($\text{QBER}\ge 9.8\%$ at $f_{EC}\!=\!1.16$) fire at approximately the same noise level: $p\!\approx\!0.17$ for CHSH and $p\!\approx\!0.11$ for QBER. The QBER-based abort is therefore slightly more conservative than the CHSH-based abort under this noise model, a finding consistent with published analyses which note that Bell-inequality violation is a \emph{sufficient} but not always \emph{tightest} indicator of security~\cite{Scarani2009, ekertQI}. In practice the two criteria are used in combination: CHSH offers a physics-level guarantee that is independent of parameter estimation assumptions, while QBER provides a tighter numerical abort threshold. Note that an Eve-interception sweep for E91 is not reported in this work; the current \texttt{EveInterceptor} is tailored to the single-qubit P\&M protocols and the entangled-pair attack model is flagged as future work in \cref{sec:future-work}.

% ─────────────────────────────────────────────────────────────────────────────
\section{Three-Protocol Noise Tolerance Benchmark}
\label{sec:results-noise-benchmark}

\Cref{fig:qber-noise-three} presents the QBER response of BB84, B92, and E91 under depolarising noise on an otherwise adversary-free channel, directly addressing O3.1 and O5.3. The common $x$-axis and matched benchmarking conditions allow the three protocols to be ranked against one another.

\begin{figure}[h]
    \centering
    \includegraphics[width=\linewidth]{images.sim/qber_vs_noise.png}
    \caption{QBER versus depolarising noise strength for BB84, B92, and E91 ($n=1000$ per trial, $N=20$). The E91 simulation uses the \texttt{both} channel topology. The horizontal dashed lines mark the Shor--Preskill thresholds for BB84 (11\%) and B92 (6.5\%); the corresponding practical $f_{EC}\!=\!1.16$ thresholds sit at 9.8\% and 6.5\% respectively.}
    \label{fig:qber-noise-three}
\end{figure}

The observed QBER curves match their analytical predictions closely. BB84's slope of $1/2$ ($\text{QBER} = p/2$) follows from the fact that depolarising noise randomises the qubit outcome with probability $p$, and a random outcome agrees with Alice's bit with probability $1/2$. B92 tracks the slightly steeper profile $\text{QBER} = p/(1+p)$ due to its conclusive-outcome filtering~\cite{B92}, implemented in \texttt{B92Protocol.theoretical\_qber()}. E91's \texttt{both}-topology curve rises with the quadratic correction $\text{QBER} = p - p^{2}/2$ because both qubits of the singlet are independently depolarised before measurement.

The threshold-crossing noise levels rank the three protocols clearly:
%
\begin{itemize}
    \item BB84 crosses the $f_{EC}\!=\!1.16$ threshold (9.8\% QBER) at $p_{\mathrm{dep}}\!\approx\!0.20$.
    \item E91 crosses the $f_{EC}\!=\!1.16$ threshold (9.8\% QBER) at $p_{\mathrm{dep}}\!\approx\!0.11$.
    \item B92 crosses the $f_{EC}\!=\!1.16$ threshold (6.5\% QBER) at $p_{\mathrm{dep}}\!\approx\!0.07$.
\end{itemize}
%
BB84 is therefore the most noise-tolerant of the three protocols by approximately a factor of two over E91 and nearly a factor of three over B92. The E91 disadvantage relative to BB84 is not a property of the entanglement-based paradigm itself but of the \texttt{both} topology chosen in this benchmark: applying depolarising noise to both halves of the singlet exposes the Bell state to compounded decoherence. In the \texttt{bob}-only topology (where only one qubit traverses a noisy channel) the analytical prediction collapses to $\text{QBER}=p/2$, matching BB84 exactly; this alternative configuration is supported by the \texttt{e91\_channel\_topology} parameter in \texttt{SimConfig} and will be added to the appendix for completeness. B92's poorer performance is a structural consequence of its encoding geometry: non-orthogonal state discrimination is unavoidably harder than orthogonal-basis discrimination, and the protocol's conclusive-outcome sifting amplifies small noise-induced errors into conclusive-round discrepancies.

% ─────────────────────────────────────────────────────────────────────────────
\section{Information-Theoretic Capacity Benchmark}
\label{sec:results-mutual-info}

QBER alone is an incomplete performance indicator because it does not account for how many sifted bits a protocol actually produces per transmitted qubit. \Cref{fig:mutual-info} resolves this by plotting the Alice--Bob mutual information, normalised to bits per 100 \emph{transmitted} qubits, as a function of depolarising noise strength. This representation, defined by the \texttt{mutual\_information} property in \texttt{QKDResult} as $\eta_{\mathrm{sift}}(1 - h(\text{QBER}))$, combines sifting efficiency and channel quality into a single information-theoretic throughput metric.

\begin{figure}[h]
    \centering
    \includegraphics[width=\linewidth]{images.sim/mutual_info_vs_noise.png}
    \caption{Alice--Bob mutual information (bits per 100 transmitted qubits) versus depolarising noise strength for the three protocols ($n=1000$, $N=20$, E91 \texttt{both} topology). The $y$-axis is scaled by sifting efficiency, so the zero-noise intercepts reproduce the theoretical sifting rates of 50, 25, and 22.2 bits per 100 qubits for BB84, B92, and E91 respectively.}
    \label{fig:mutual-info}
\end{figure}

At zero noise the three curves meet their analytical intercepts. BB84 sits at $\sim\!50$ bits per 100 transmitted qubits, because each retained sifted bit carries exactly one bit of Alice--Bob mutual information ($1 - h(\text{QBER}) = 1$) and the sifting efficiency is 50\%. B92 and E91 sit at $\sim\!25$ and $\sim\!22$ respectively, set by their $25\%$ and $22.2\%$ sifting efficiencies. The horizontal separation between the three curves is therefore structural: it reflects the sifting geometry of each protocol, not a noise-tolerance difference.

The vertical spread is noise-driven. Each curve decays as $\eta_{\mathrm{sift}} \cdot (1 - h(\text{QBER}(p)))$, where $\eta_{\mathrm{sift}}$ is the zero-noise intercept and $\text{QBER}(p)$ is the protocol-specific noise response derived in \cref{sec:results-noise-benchmark}. BB84 retains $\sim\!19$ bits per 100 qubits at $p = 0.30$, compared to $\sim\!8$ bits for B92 and $\sim\!4$ bits for E91 (\texttt{both}-topology). BB84's structural advantage is therefore amplified, not reduced, by channel noise: the high-sifting-efficiency protocol also decays more slowly.

This benchmark addresses objective O4.2 by providing a single-axis performance metric for inter-protocol comparison. The immediate engineering implication is that for any noise level below the B92 security threshold, BB84 generates sifted key material at roughly $2.5\times$ the rate of B92 (19 vs 8 bits per 100 qubits at $p=0.30$) and roughly $5\times$ the rate of E91 (\texttt{both}-topology). For throughput-limited applications --- high-bandwidth sector key refresh, for instance --- this ratio is a dominant consideration and would typically tip protocol selection toward BB84 unless E91's device-independence or physics-level detection argument is required by the deployment threat model.

% ─────────────────────────────────────────────────────────────────────────────
\section{Secure Key Rate and Deployment Feasibility}
\label{sec:results-keyrate}

\textit{[Sign-post: the key-rate-versus-QBER and key-rate-versus-distance figures are pending and will be inserted here. The discussion below establishes the engineering framing; once the figures are generated the commentary will be extended to reference specific maximum-distance values.]}

For prepare-and-measure QKD over fibre, the GLLP formula of \cref{eq:gllp} translates the QBER benchmarks of \cref{sec:results-noise-benchmark} directly into secure-key fractions. With $f_{EC}\!=\!1.16$, BB84 sustains a positive $\keyrate$ up to 9.8\% QBER and B92 up to 6.5\% QBER; above these thresholds no classical post-processing can extract a secret key. Using the standard fibre attenuation coefficient $\alpha\!=\!0.2$~dB/km and a dark-count-limited QBER floor, these thresholds imply maximum deployment distances that distinguish the two protocols even when their baseline channel quality is identical. BB84's higher threshold allows roughly $\log(9.8/6.5)/0.2\!\approx\!9$~km of additional tolerable attenuation before the link aborts, which is comparable to the metropolitan-ring scale used in recent SDN-QKD demonstrations~\cite{Zhang2025}. The gap between $f_{EC}\!=\!1.0$ and $f_{EC}\!=\!1.16$ thresholds --- an often-overlooked $\sim\!1.2$ percentage points of tolerable QBER for BB84 --- maps directly to a further several kilometres of reach and is of direct engineering significance for deployment studies.

% ─────────────────────────────────────────────────────────────────────────────
\section{Limitations and Sources of Error}
\label{sec:limitations-results}

The results above are constrained by the scope boundaries set in \cref{sec:modelling-assumptions}. Four sources of uncertainty bear on the quantitative statements in this chapter.

\textit{Statistical sampling.} Each data point is averaged over $N\!=\!20$ trials of $n\!=\!1000$ qubits, giving a standard error of order $\sigma/\sqrt{N}$; this is visible as the small error bars in every figure. Threshold-crossing points reported to two decimal places (e.g.\ $p\!\approx\!0.16$) therefore carry an uncertainty of order $\pm 0.02$. Increasing $N$ or $n$ would tighten these bounds at a linear cost in simulation time.

\textit{Single-photon source assumption.} The simulator assumes one photon per round; real weak-coherent-pulse sources have a non-zero probability of emitting two or more photons, enabling the photon-number-splitting attack. Decoy-state BB84~\cite{Lo2005Decoy, Ma2005} is the standard practical countermeasure and is the natural next extension of the platform.

\textit{Trusted devices.} All preparation and measurement operations are assumed noiseless; detector-efficiency mismatch, timing-jitter, and afterpulsing are excluded. These side-channel imperfections are the primary attack surface in real deployments~\cite{Scarani2009} and are outside the scope of this platform.

\textit{Classical fibre attenuation.} Photon loss is modelled as a classical post-selection rather than a full quantum channel. The effect on transmission efficiency is captured correctly but the quantum-optical description of loss (which affects decoy-state analysis and finite-key bounds) is not reproduced.

These are deliberate scope choices, not implementation defects. Acknowledging them precisely is important because the protocol-ranking conclusions of this chapter are valid within this scope and should not be extrapolated to, for example, weak-coherent-pulse BB84 against a device-independent adversary without extending the noise model accordingly.
% >>> END ADDED SECTION <<<



% >>> ADDED SECTION: Conclusion & Future Work Chapter <<<
% ═══════════════════════════════════════════════════════════════════════════════
\chapter{Conclusion \& Future Work}
% ═══════════════════════════════════════════════════════════════════════════════

% ─────────────────────────────────────────────────────────────────────────────
%  CONCLUSION -- Section I: Achievements Against Objectives
% ─────────────────────────────────────────────────────────────────────────────

\section{Achievements Against Objectives}

This project set out to design, implement, and evaluate a modular simulation platform for comparing discrete-variable QKD protocols under realistic conditions. The platform successfully implements all three target protocols: BB84, B92, and E91.

Objective group O1 (Literature Review and Analytical Framework) was achieved through a comprehensive review spanning quantum mechanical foundations, protocol theory, post-processing pipelines, noise models, and simulation frameworks. The review established the parameter choices and metrics used throughout the project.

Objective group O2 (Protocol Implementation) was fully achieved. BB84 was validated against a sifting efficiency target of $50\% \pm 2\%$ and a full intercept-resend QBER of $25\% \pm 2\%$. B92 was validated at $25\% \pm 2\%$ sifting efficiency. E91 was implemented with CHSH parameter convergence to $2\sqrt{2}$ in the absence of eavesdropping.

Objective group O3 (Channel Modelling and Eavesdropper Simulation) was achieved through the implementation of depolarising, phase damping, and amplitude damping noise models applied via the channel-marker approach, along with a configurable intercept-resend eavesdropper.

Objective group O4 (Post-Processing and Key Rate Analysis) was achieved through GLLP key rate computation with $f_{EC} = 1.16$ and mutual information analysis for both BB84 and B92. The security thresholds of $\sim$9.8\% (BB84) and $\sim$6.5\% (B92) were confirmed.

Objective group O5 (Verification and Benchmarking) was largely achieved. All ideal-channel and eavesdropper validation targets were met, and a comprehensive visualisation suite was produced. Quantitative comparison against published experimental data (O5.2) was performed at a high level, with discrepancies attributed to the modelling assumptions documented in the methodology.

Objective group O6 (Real-World Feasibility) was addressed through the key-rate-versus-distance analysis and the broader context discussion that follows.

% NOTE: Update pass/fail status and numerical values above based on your actual results.


% ─────────────────────────────────────────────────────────────────────────────
%  CONCLUSION -- Section II: Key Findings
% ─────────────────────────────────────────────────────────────────────────────

\section{Key Findings}

Three findings from this project merit particular emphasis. First, the distinction between $f_{EC} = 1.0$ and $f_{EC} = 1.16$ has real engineering consequences: the BB84 security threshold shifts from $\sim$11\% to $\sim$9.8\% QBER, a reduction that directly affects maximum deployment distance calculations. Simulation studies that assume Shannon-limit error correction overestimate practical system performance.

Second, phase damping noise reveals protocol-specific vulnerabilities that are invisible under depolarising noise alone. B92's asymmetric two-state encoding ($\ket{0}$, $\ket{+}$) creates a bit-value-dependent error rate under phase noise, a behaviour not exhibited by BB84's symmetric four-state scheme. This finding demonstrates the value of testing protocols under multiple noise models rather than relying solely on the depolarising channel.

Third, E91's CHSH-based eavesdropper detection provides qualitatively different security evidence from the QBER-based detection used by prepare-and-measure protocols. The entanglement-based approach detects eavesdropping through the degradation of a physical property (Bell inequality violation) rather than through statistical error estimation, offering a conceptually distinct and complementary security mechanism.


% ─────────────────────────────────────────────────────────────────────────────
%  CONCLUSION -- Section III: Broader Context
% ─────────────────────────────────────────────────────────────────────────────

\section{Broader Context}

QKD and Post-Quantum Cryptography are complementary rather than competing approaches to quantum-safe security. PQC provides computational security based on mathematical problems believed to resist quantum attack, while QKD provides information-theoretic security grounded in the laws of physics. A defence-in-depth strategy would deploy both: PQC to protect the vast majority of communications at low cost, and QKD to secure the most sensitive channels where long-term confidentiality is essential.

The "harvest now, decrypt later" threat model gives QKD deployment an urgency that extends beyond the timeline for building a cryptographically relevant quantum computer. Adversaries intercepting and storing encrypted traffic today may be able to decrypt it in the future once quantum computing capabilities mature. For sectors handling data with long confidentiality lifetimes, including government and defence, financial services, healthcare, and critical national infrastructure, the window of vulnerability is already open.

However, QKD deployment faces significant economic and infrastructure barriers. Current QKD hardware is expensive, fibre infrastructure is geographically limited, and the trusted-node problem constrains the distances over which key distribution can be performed without intermediate relay points. These constraints mean that near-term QKD deployment is most likely in high-value, point-to-point links rather than as a wholesale replacement for classical key exchange.

The simulation platform developed in this project contributes to this landscape by enabling protocol selection and feasibility analysis without requiring access to expensive quantum hardware. By providing an accessible tool for comparing BB84, B92, and E91 under realistic conditions, the platform lowers the barrier to informed deployment decisions for organisations evaluating QKD as part of their quantum-safe migration strategy.


% ─────────────────────────────────────────────────────────────────────────────
%  CONCLUSION -- Section IV: Project Management Reflection
% ─────────────────────────────────────────────────────────────────────────────

\section{Project Management Reflection}

The project timeline spanned approximately eight weeks from initial literature review through final report submission. The work was structured in phases: literature review and scope definition in weeks one and two, BB84 implementation and validation in weeks three and four, B92 implementation and comparative benchmarking in week five, E91 implementation in week six, and report writing with final analysis in weeks seven and eight.

Several aspects of the project proceeded as planned. The BB84 implementation and validation were completed on schedule, and the literature review scope was defined early enough to avoid scope creep during the writing phase. The decision to adopt YAML configuration files mid-project proved valuable for reproducibility, as it allowed all simulation parameters to be version-controlled alongside the code.

The most significant deviation from the original plan concerned E91. It was initially designated as a stretch goal, contingent on the progress of BB84 and B92 development. The modular architecture of the codebase, particularly the separation of noise and benchmarking modules from protocol-specific logic, ultimately enabled E91 to be integrated more efficiently than anticipated. A mid-project refactoring from a monolithic BB84 script into the modular package structure required short-term effort but substantially reduced the cost of adding subsequent protocols.

In terms of personal development, this project provided practical experience with quantum information theory and the Qiskit framework, as well as exposure to the challenges of translating theoretical physics into validated software. The iterative development approach, where each protocol was validated against theoretical predictions before proceeding to the next, proved effective for maintaining confidence in the simulation outputs throughout the project.


% ─────────────────────────────────────────────────────────────────────────────
%  CONCLUSION -- Section V: Limitations
% ─────────────────────────────────────────────────────────────────────────────

\section{Limitations}

The limitations of this work are discussed in detail in \cref{sec:limitations-results}. In summary, the simulation operates under single-photon source, trusted-device, and authenticated classical channel assumptions that exclude several classes of practical imperfection. These are deliberate scope boundaries that maintain focus on protocol-level comparison within the constraints of an undergraduate project.


% ─────────────────────────────────────────────────────────────────────────────
%  CONCLUSION -- Section VI: Future Work
% ─────────────────────────────────────────────────────────────────────────────


% potentially remove/reduce definitively too long
\section{Ethical \& EDI Considerations}

This project is a software-based simulation study with no human participants, personal data collection, or physical experimentation, and therefore required no formal ethical approval. The simulation platform and all associated code are the original work of the author, developed using open-source tools (Python, Qiskit) that are freely available to researchers regardless of institutional affiliation or geographic location. The choice of an open-source, Python-based implementation was made in part to ensure that the platform is accessible to a broad range of users, including those in resource-constrained research environments. No datasets containing personal or sensitive information were used at any stage of the project. The broader implications of QKD technology, including questions of equitable access to quantum-safe security infrastructure, are discussed in the Broader Context section above.


% review and reduce
\section{Future Work}
\label{sec:future-work}

The simulation platform developed in this project provides a foundation that can be extended in several directions.

\begin{itemize}
    \item \textbf{Decoy-State BB84:} Implementing the decoy-state protocol \cite{Lo2005Decoy} would relax the single-photon source assumption, enabling simulation of weak coherent pulse sources and photon-number-splitting attacks. % TODO: Add detail on expected impact on key rate curves.

    \item \textbf{Continuous-Variable QKD:} Integrating CV protocols such as GG02 would enable direct comparison between discrete-variable and continuous-variable approaches on the same benchmarking platform. % TODO: Add detail on homodyne detection modelling.

    \item \textbf{Network-Scale Simulation:} Extending the platform to model multi-node QKD networks using a discrete-event network simulator would enable analysis of trusted-node routing and network topology effects. % TODO: Add detail on specific network scenarios.

    \item \textbf{Experimental Validation:} Quantitative comparison against published experimental data, such as the results of Boaron \textit{et al.} \cite{Boaron2018}, would identify where modelling assumptions introduce the largest discrepancies. % TODO: Add detail on specific datasets for comparison.

    \item \textbf{Advanced Attack Models:} Incorporating collective and coherent attack strategies beyond the intercept-resend model would test protocol security under more sophisticated adversarial assumptions. % TODO: Add detail on specific attack types.

    \item \textbf{Machine Learning Integration:} Applying ML techniques for adaptive protocol parameter optimisation and anomaly-based eavesdropper detection could improve the platform's practical utility. % TODO: Add detail on specific ML approaches.
\end{itemize}



% ─────────────────────────────────────────────────────────────────────────────
% REFERENCES
% ─────────────────────────────────────────────────────────────────────────────
% IEEE references: 8pt, 9pt line spacing, bracketed Arabic numerals [1], [2]
\clearpage
{\footnotesize
\bibliographystyle{IEEEtran}
\bibliography{bib}
}

% ─────────────────────────────────────────────────────────────────────────────
% APPENDICES
% ─────────────────────────────────────────────────────────────────────────────
\appendix

\chapter{Source Code}
% TODO: Full source code listings or link to repository.

\chapter{Additional Results}
% TODO: Supplementary plots, extended data tables.

\end{document}
